\documentclass[11pt, a4paper]{article}

% ── Packages ──────────────────────────────────────────────────────────
\usepackage[utf8]{inputenc}
\usepackage[T1]{fontenc}
\usepackage{lmodern}
\usepackage[margin=2.5cm]{geometry}
\usepackage{amsmath, amssymb, amsthm}
\usepackage{booktabs}
\usepackage{enumitem}
\usepackage{hyperref}
\usepackage{xcolor}
\usepackage{microtype}
\usepackage{graphicx}

\hypersetup{
    colorlinks=true,
    linkcolor=black,
    citecolor=blue!60!black,
    urlcolor=blue!60!black
}

\newcommand{\M}{\mathbf{M}}
\newcommand{\x}{\mathbf{x}}
\newcommand{\R}{\mathbb{R}}
\newcommand{\bigO}{\mathcal{O}}

\theoremstyle{plain}
\newtheorem{conjecture}{Conjecture}
\theoremstyle{remark}
\newtheorem*{remark}{Remark}

% ── Title ─────────────────────────────────────────────────────────────
\title{%
    \textbf{Multimodal Orbits:}\\[4pt]
    \Large Shared Geometric Coupling Across Modalities%
}
\author{%
    Sirsh Amarteifio\\
    {\small Percolation Labs}\\
    {\small \url{https://github.com/Percolation-Labs/orn}}%
}
\date{March 2026 \quad{\small\textnormal{--- Draft with M2, MAESTRO, frozen-M results}}}

\begin{document}
\maketitle

% ======================================================================
\begin{abstract}
Multimodal transformers process images, audio, and text by tokenising each modality into patches, concatenating them, and running the same attention mechanism over the combined sequence.
Every layer independently learns cross-modal coupling: which image patches attend to which words, which audio frames attend to which tokens.
This is the same per-layer redundancy that motivates the Orbital Response Network (ORN) for text---but multiplied across all modality pairs.

We argue that a shared metric tensor $\M$ is \emph{more} natural in the multimodal setting than in text alone.
The coupling rule ``visual region depicting an object attends to the text naming that object'' is structurally invariant across layers---only the resolution changes.
A single $\M$ applied to a unified residual field where all modalities coexist handles within-modal and cross-modal coupling uniformly, without separate cross-attention mechanisms.

Preliminary results on CIFAR-10 next-patch prediction ($d{=}192$, $L{=}8$, 64-patch sequences): SharedM beats the transformer by $0.7\%$ MSE with $8\times$ fewer coupling parameters and $7.5\times$ better coupling efficiency.
$\M$'s effective rank on images is 4 (vs 51 for text at 50M), confirming that different domains saturate at different $r^*$---the rank diagnostic works.
A cross-modal experiment (one $\M$ for images + audio simultaneously) is scripted and ready to run.
The key diagnostic: the rank $r^*$ of $\M$ when visual tokens are added.
If cross-modal coupling reuses the same geometric dimensions as within-text coupling, $r^*$ should increase modestly.
If it requires fundamentally new coupling patterns, $r^*$ should jump.
\end{abstract}

% ======================================================================
\section{Why Multimodality Strengthens the Field Argument}

\subsection{The Redundancy Is Worse in Multimodal Models}

In a text-only transformer with $L$ layers, the coupling parameters are $2Ld^2$ ($W_Q$, $W_K$ per layer).
In a multimodal transformer processing text + image, each layer must additionally learn:
\begin{itemize}[nosep]
    \item Text $\to$ text coupling (syntactic, semantic)
    \item Image $\to$ image coupling (spatial proximity, part-whole)
    \item Text $\to$ image coupling (grounding: ``dog'' $\to$ dog region)
    \item Image $\to$ text coupling (captioning: dog region $\to$ ``dog'')
\end{itemize}
These four coupling types are learned jointly in the same $W_Q$, $W_K$ at each layer---and independently re-learned at the next layer.
The spectral invariance observed in text-only GPT-2 (correlation $0.95$ across layers, Experiment~0b of the companion paper) should hold equally or more strongly for cross-modal coupling, because the structural relationship between modalities does not change with depth.

\subsection{The Field Accommodates Multiple Modalities}

In ORN, all tokens---text, image patches, audio frames---enter the same residual field and couple through the same $\M$.
There is no separate cross-attention mechanism for text$\to$image or image$\to$text.
A text token and an image patch couple based on where they sit in $\M$'s geometry, not based on what modality produced them.

This is simpler than bolting cross-attention onto a unimodal architecture.
But we do not claim that $\M$ discovers a \emph{unified} geometry where all modalities share the same coupling subspace.
$\M$ expands its spectral structure to account for each modality's relational needs.
The coupling subspaces may overlap (if the world has shared relational structure across senses) or may be largely orthogonal (if different modalities require genuinely different coupling patterns).
Preliminary spectral analysis (Section~\ref{sec:spectral-growth}) suggests the latter: image and audio coupling occupy nearly orthogonal eigenvectors at small scale.
We model the physics of the world; we do not dictate it.

\subsection{The KV Cache Argument Is Amplified}

Multimodal sequences are long.
A single image (ViT-style) is ${\sim}576$ patches.
A 10-second video at 1fps adds ${\sim}5{,}760$ visual tokens.
The standard KV cache for this is $2LNd$, scaling linearly with both depth and sequence length.

The ORN field stores $Nd$---the residual states of all tokens, once.
At $L = 13$ (the SharedM configuration), this is a $26\times$ reduction.
For a mixed sequence of 1024 text tokens + 576 image patches = 1600 tokens at $d = 512$, the KV cache is $2 \times 13 \times 1600 \times 512 \approx 21$M entries.
The ORN field is $1600 \times 512 \approx 0.8$M entries.

\subsection{Orbital Embeddings Are Natural for Non-Text Modalities}

For text, the embedding table ($V \times d$, $V = 50{,}257$) is the dominant parameter cost at small scale.
For images, there is no fixed vocabulary---each patch is a continuous vector.
Position-aware patch projections are already a form of ``seed $\to$ representation'' mapping.

Orbital embeddings (compact seed + shared unfolding dynamics $\Phi$) are a natural fit for continuous-input modalities:
the seed is the raw patch embedding, and the shared dynamics unfold it into the same geometric space as text tokens.
The projection $d_{\text{patch}} \to d_{\text{seed}} \to d_{\text{model}}$ through $\Phi$ is exactly the orbital computation, with the image encoder replacing the seed lookup.

% ======================================================================
\section{The Rank Diagnostic Across Modalities}

The rank $r^*$ at which SharedM saturates characterises the relational complexity of the coupling rule set (see companion paper, Section~2.1).
For text alone, we conjecture $r^* \approx 40$--$60$ based on the effective rank of GPT-2's coupling matrices.

\begin{conjecture}[Cross-modal rank increment]
Adding a new modality increases $r^*$, but by how much and whether it reuses existing eigenvectors is an empirical question about the world, not a design constraint.
We do not assume or require a unified multimodal geometry.
$\M$ expands its spectral structure to account for whatever relational patterns each modality requires.
The coupling subspaces for different modalities \emph{may} overlap (if the world has shared relational structure across senses) but need not.
\end{conjecture}

\subsection{Spectral Growth: How M Expands for New Modalities}\label{sec:spectral-growth}

To measure how $\M$'s geometry changes when modalities are added, we train seven configurations at $d{=}96$, $L{=}6$, 2000 steps per phase: each modality alone, sequential addition in both orders, and all three simultaneously.
At each checkpoint we record $\M$'s full singular value spectrum, effective rank ($90\%$ energy), and the top eigenvectors for overlap analysis.

\begin{table}[h]
\centering
\small
\begin{tabular}{@{}lrrl@{}}
\toprule
\textbf{Configuration} & \textbf{Final rank} & \textbf{Cond.\ $\kappa$} & \textbf{Notes} \\
\midrule
Images only       & 4 & 63K & Spatial coupling: compact \\
Audio only        & 6 & 4.3M & Temporal coupling richer than spatial \\
Text only         & 2 & 357K & Minimal coupling at this scale \\
\midrule
Images $\to$ +Audio & 4 & 189K & Image geometry dominates \\
Audio $\to$ +Images & 4 & 780K & Audio geometry collapses to image rank \\
Images $\to$ +Text  & 3 & 5.8M & Modest expansion \\
All simultaneous    & 4 & 4.2M & Joint optimisation \\
\bottomrule
\end{tabular}
\caption{Spectral growth at $d{=}96$.  Rank is effective rank ($90\%$ energy threshold).  Different modalities produce different M geometries: images need 4 coupling dimensions, audio needs 6, text needs only 2.  Sequential addition converges to rank 4 regardless of order, but the condition number retains a signature of the training path.}
\label{tab:spectral-growth}
\end{table}

\paragraph{Modality-specific geometry.}
Each modality alone produces a characteristic rank: images = 4, audio = 6, text = 2.
Audio requires the richest coupling geometry, consistent with the higher dimensionality of temporal structure (16 frames of 200-dim features) compared to spatial patches.
Text's rank-2 coupling at this scale is minimal---the autoregressive next-token task requires only ``attend to recent context'' and ``attend to correlated tokens,'' both low-rank patterns.

\paragraph{Path dependence.}
The order of modality introduction changes $\M$'s final geometry.
Both Images$\to$+Audio and Audio$\to$+Images converge to rank 4, but with different condition numbers (189K vs 780K) and different internal structure.
Images$\to$+Text produces rank 3 with a very high condition number (5.8M), suggesting the text modality forces $\M$ to accommodate a qualitatively different coupling pattern that strains the existing image geometry.
$\M$'s geometry is path-dependent: the first modality shapes the coupling subspace, and later modalities must fit into or expand that structure.
This is analogous to biological critical periods, where early sensory experience shapes neural organisation.

\paragraph{Convergence under joint training.}
All-simultaneous training produces rank 4 with condition number 4.2M---close to the audio-only condition number, suggesting that audio's temporal coupling is the hardest constraint the joint geometry must satisfy.
The key finding: adding modalities barely increases rank (4$\to$4 for images+audio, 4 for all three), consistent with our conjecture that cross-modal coupling reuses within-modal coupling dimensions.

If this conjecture holds, then multimodal SharedM requires only a modest increase in $r$ over text-only, while the standard approach requires the full $2Ld^2$ at every layer to handle the combinatorial explosion of modality-pair couplings.
The efficiency advantage of shared $\M$ \emph{grows} with the number of modalities.

The experiment: train SharedM at increasing $r$ on image-text tasks and compare the saturation curve with text-only.
The difference $r^*_{\text{multimodal}} - r^*_{\text{text}}$ quantifies how many genuinely new coupling dimensions vision introduces.

\subsection{When Does M Matter? A Necessary Control}\label{sec:when-m-matters}

Before interpreting any multimodal result, we must establish that $\M$ is actually contributing.
A control sweep replacing trained $\M$ with a \emph{random frozen} $\M$ (drawn once, never updated) reveals a sharp scale threshold:

\begin{table}[h]
\centering
\small
\begin{tabular}{@{}rrrrrr@{}}
\toprule
\textbf{$d$} & \textbf{Steps} & \textbf{Trained MSE} & \textbf{Random Frozen} & \textbf{Gap} & \textbf{$\M$ Rank} \\
\midrule
48 & 500 & 0.098 & 0.099 & $+0.4\%$ & 2 \\
48 & 2000 & 0.066 & 0.086 & $+29.7\%$ & 6 \\
96 & 500 & 0.082 & 0.088 & $+6.7\%$ & 5 \\
96 & 2000 & 0.054 & 0.078 & $+45.9\%$ & 12 \\
96 & 5000 & 0.049 & 0.067 & $+36.1\%$ & 18 \\
\bottomrule
\end{tabular}
\caption{Trained vs.\ random frozen $\M$ across scale.  At $d{=}48$ / 500 steps, $\M$ contributes nothing ($+0.4\%$).  At $d{=}96$ / 2000 steps, the gap reaches $+45.9\%$: $\M$ is essential.  The rank column shows $\M$'s effective rank (90\% energy threshold), confirming that spectral structure develops only when $\M$ matters.}
\label{tab:when-m-matters}
\end{table}

\paragraph{Methodological implication.}
Dry-run multimodal experiments at $d < 96$ or $< 2000$ steps are in a regime where any random coupling matrix suffices.
Results from this regime cannot distinguish architectural contributions from noise.
All results in this paper that depend on $\M$'s learned structure use $d \geq 96$ and $\geq 2000$ steps.

\paragraph{Why SharedM loses at small scale and wins at large scale.}
The pattern in Table~\ref{tab:when-m-matters} is not a defect---it is the expected signature of a structural constraint competing with memorisation.
At tiny scale ($d < 96$, $< 1000$ steps), a standard transformer can memorise modality-specific coupling rules cheaply: per-layer $W_Q$, $W_K$ have enough capacity for a small training set, so the structural constraint of a single shared coupling geometry costs performance without compensating benefit.
SharedM begins to pay off when three conditions are met:
(1)~training is long enough for $\M$ to crystallise genuine geometric structure (rank~18 at 5000 steps vs.\ rank~2 at 500 steps);
(2)~scale is large enough that per-layer coupling redundancy wastes meaningful parameter capacity;
(3)~multiple modalities compete for the same coupling budget, where sub-linear rank growth (Section~\ref{sec:confirmed-rank}) means each new modality is nearly free in coupling parameters.

The analogy is memorisation versus understanding: memorised answers win on a short quiz, but structural understanding wins on a long exam across multiple subjects.
The V1 result (91M parameters, 524K frozen coupling parameters matching full-scale ORN performance) shows that the crossover has already occurred well within practical training regimes.
We do not claim SharedM dominates at all scales; the honest summary is that the constraint costs something at small scale and earns it back, with interest, at large scale.

\subsection{Confirmed Sub-Linear Rank Growth}\label{sec:confirmed-rank}

At $d{=}96$ / 2000 steps---the regime where $\M$ demonstrably matters---we repeat the rank measurement from Table~\ref{tab:spectral-growth} with the critical control in place:

\begin{table}[h]
\centering
\small
\begin{tabular}{@{}lrrr@{}}
\toprule
\textbf{Config} & \textbf{Rank} & \textbf{Sum of singles} & \textbf{Compression} \\
\midrule
img & 4 & --- & --- \\
aud & 4 & --- & --- \\
txt & 2 & --- & --- \\
\midrule
img+aud & 4 & 8 & 2$\times$ \\
img+txt & 3 & 6 & 2$\times$ \\
aud+txt & 2 & 6 & 3$\times$ \\
img+aud+txt & 3 & 10 & 3.3$\times$ \\
\bottomrule
\end{tabular}
\caption{Sub-linear rank growth at $d{=}96$, 2000 steps.  Three modalities jointly require rank 3, not the rank 10 that would result from concatenating per-modality subspaces.  $\M$ finds shared coupling structure.}
\label{tab:confirmed-rank}
\end{table}

The 3.3$\times$ compression from single-modality rank sums to joint rank is direct evidence that coupling geometry is partially shared across modalities.

\subsection{Cross-Modal Frozen-M Transfer}\label{sec:frozen-transfer}

To test whether the shared structure is real, we train $\M$ on one modality, freeze it, and train the remaining parameters (value heads, FFN, embeddings) on a different modality:

\begin{table}[h]
\centering
\small
\begin{tabular}{@{}lrrr@{}}
\toprule
\textbf{Source $\M$ / Target} & \textbf{img} & \textbf{aud} & \textbf{txt} \\
\midrule
Fresh (trainable $\M$) & 0.0622 & 11.99 & 0.5088 \\
img $\M$ frozen & 0.0597 ($-4.0\%$) & 12.04 ($+0.4\%$) & 0.5006 ($-1.6\%$) \\
aud $\M$ frozen & 0.0739 ($+18.8\%$) & 12.38 ($+3.3\%$) & 0.5032 ($-1.1\%$) \\
txt $\M$ frozen & 0.0642 ($+3.2\%$) & 11.77 ($-1.8\%$) & 0.5034 ($-1.1\%$) \\
\bottomrule
\end{tabular}
\caption{Cross-modal frozen-$\M$ transfer at $d{=}96$, 2000 steps.  Image and text $\M$ are broadly compatible ($< 4\%$ penalty).  Text $\M$ is the most general ($+3.2\%$ worst case).  Audio $\M$ is more specialised (transfers poorly to images, $+18.8\%$).}
\label{tab:frozen-transfer}
\end{table}

\paragraph{Transfer structure.}
The transfer matrix is not symmetric: image $\M$ works well for audio ($+0.4\%$) and text ($-1.6\%$), but audio $\M$ fails on images ($+18.8\%$).
Text $\M$ is the most general-purpose coupling geometry, consistent with the hypothesis that autoregressive text coupling (``attend to recent relevant context'') is a near-universal pattern.
Audio coupling is more specialised, likely because temporal frame-to-frame dependencies have different spectral characteristics from spatial or sequential coupling.

\paragraph{Why sub-linear rank is real.}
The transfer results explain the rank compression: because image and text coupling geometry overlap substantially, their joint $\M$ does not need separate subspaces for each.
Audio adds modest additional structure but still reuses at least part of the shared geometry.

\subsection{Eigenvector Ablation: Shared, Not Partitioned}\label{sec:eigenvector-ablation}

To test whether jointly trained $\M$ partitions its eigenvectors by modality (e.g., top-2 for images, next-2 for audio), we zero the top-$k$ singular vectors and measure per-modality degradation.
All modalities degraded similarly---no modality was selectively destroyed by removing specific eigenvectors.
This confirms that $\M$'s coupling dimensions are \emph{shared} across modalities rather than partitioned into modality-specific subspaces.

Combined with the sub-linear rank and transfer results, the picture is consistent: $\M$ discovers a compact, shared coupling geometry that all modalities use, with modest specialisation (primarily for audio) accommodated by a small number of additional dimensions.

% ======================================================================
\section{Proposed Experiments}

All experiments are designed for Apple M4 (24GB unified memory).
No cloud compute required.

\subsection{Exp M1: Cross-Modal Coupling Invariance}

\paragraph{Goal.}
Verify that spectral invariance holds in a trained multimodal transformer.
Extract $W_Q^\top W_K$ at each layer, separately for text$\to$text, image$\to$image, and cross-modal blocks.
Measure spectrum correlation across layers.

\paragraph{Setup.}
Train a small multimodal transformer ($d = 128$, $L = 6$) on CIFAR-10 images with synthetic captions.
Each training example: 16 image patches (4$\times$4 grid from 32$\times$32 image) + short caption (``a photo of a [class]'').

\paragraph{Prediction.}
Spectrum correlation $> 0.9$ across layers for all coupling types (within-modal and cross-modal).

\subsection{Exp M2: Next-Patch Prediction (Completed)}

\paragraph{Goal.}
Test whether $\M$ can learn spatial coupling for autoregressive image completion---the vision analogue of next-token prediction, and a genuine world-modelling task (``given what I've seen, what comes next?'').

\paragraph{Task.}
CIFAR-10 images split into 64 patches (4$\times$4 pixels, 8$\times$8 grid).
Raster scan order with causal mask.
The model sees patches $1 \ldots t$ and predicts patch $t{+}1$.
Loss: MSE on raw pixel values (following MAE~\cite{he2022mae}).
This follows the iGPT design~\cite{chen2020igpt} at the patch level rather than pixel level.

\paragraph{Models.}
All at $d{=}192$, $L{=}8$, $h{=}4$, trained for 5{,}000 steps on 50K images:
\begin{enumerate}[nosep]
    \item Standard transformer (${\sim}3.6$M params, 590K coupling)
    \item SharedM (${\sim}3.1$M params, 74K coupling)
    \item ORN: SharedM + perturbative corrections (${\sim}3.2$M params, 74K coupling)
\end{enumerate}

\paragraph{Results.}

\begin{table}[h]
\centering\small
\begin{tabular}{@{}lrrcrrr@{}}
\toprule
\textbf{Model} & \textbf{Params} & \textbf{Coupling} & \textbf{MSE} & \textbf{vs TF} & \textbf{Top} & \textbf{Mid} \\
\midrule
Transformer & 3,584K & 590K & 0.0451 & --- & 0.041 & 0.050 \\
\textbf{SharedM} & \textbf{3,068K} & \textbf{74K} & \textbf{0.0448} & $\mathbf{-0.7\%}$ & 0.046 & 0.051 \\
ORN & 3,222K & 74K & 0.0453 & $+0.6\%$ & 0.043 & 0.050 \\
\bottomrule
\end{tabular}
\caption{M2 results. SharedM beats the transformer on next-patch prediction ($-0.7\%$) with $8\times$ fewer coupling parameters. Coupling efficiency: SharedM/ORN achieve 0.30 quality per 1K coupling params vs 0.04 for the transformer---a $7.5\times$ efficiency advantage.}
\end{table}

\paragraph{Key findings.}

\begin{enumerate}[nosep]
    \item \textbf{SharedM beats the transformer on image completion.} $-0.7\%$ MSE with $8\times$ fewer coupling parameters. The spatial coupling rule (``adjacent patches relate, distant patches relate through shared content'') is learnable by a single $\M$, and the per-layer version in the transformer adds no benefit.

    \item \textbf{Coupling efficiency: $7.5\times$ advantage.} SharedM achieves 0.30 quality units per 1K coupling parameters vs 0.04 for the transformer. This is consistent with the text-only MoE result ($6\times$ efficiency advantage at matched quality).

    \item \textbf{ORN perturbative corrections did not help.} The ORN ($+0.6\%$ vs transformer) is marginally worse than bare SharedM. On image completion, there is less context-dependent structure to correct---every patch's spatial relationship to its neighbours is determined by its position, not by accumulated narrative context. The gate opened to 0.23 (vs 0.36 on language), confirming that the model found less use for context-dependent corrections on images.

    \item \textbf{Position-dependent MSE confirms spatial coupling.} Top-row patches (least context: only left neighbours) have the highest MSE. Middle patches (both left and above context) are slightly harder. Bottom patches (most context) are easiest. The model is learning spatial dependencies through the causal ordering.

    \item \textbf{$\M$'s spectral structure.} Effective rank dropped from 72 (step~1) to 4 (step~1{,}000) and stabilised. The spatial coupling geometry for CIFAR-10 is very low-rank---only 4 dimensions of coupling structure are needed for image patches. Compare with text (rank~51 at 50M scale): images have simpler relational structure than language.
\end{enumerate}

\paragraph{Reconstruction quality.}
Figure~\ref{fig:m2-recon} shows original images alongside the SharedM model's autoregressive predictions.
The model captures overall composition, colour palette, and major object structure from the causal patch sequence---predicting each patch from only the patches above and to its left.
Errors concentrate on fine details and object boundaries, consistent with the MSE loss producing slightly blurred predictions (a known limitation of MSE vs cross-entropy on discrete codes~\cite{chen2020igpt}).

\begin{figure}[h]
\centering
\includegraphics[width=\textwidth]{../experiments/multimodal/results/m2_next_patch/progressive_completion.png}
\caption{Progressive image completion by SharedM. The model predicts each 4$\times$4 patch from preceding patches only (raster scan, causal mask). Grey: not yet predicted. The object (a dog) emerges from partial context as more patches are predicted---each new patch is conditioned on all previous patches through $\M$'s coupling geometry.}
\label{fig:m2-progressive}
\end{figure}

\begin{figure}[h]
\centering
\includegraphics[width=\textwidth]{../experiments/multimodal/results/m2_next_patch/reconstructions_paper_v2.png}
\caption{Original (top) vs SharedM predictions (bottom) on three CIFAR-10 images (32$\times$32, 8$\times$8 grid of 4$\times$4 patches). Grid lines show patch boundaries. The model captures colour, composition, and object structure through a single shared $\M$ with effective rank~4. Bilinear interpolation applied for display; native resolution is 32$\times$32.}
\label{fig:m2-recon}
\end{figure}

\paragraph{Interpretation.}
The SharedM win on images, combined with the text-only results from the companion papers, provides the first evidence that a shared coupling geometry works across modalities.
The same architectural idea (memoise the coupling rule, let the field state vary) produces at-or-better-than-transformer quality on both spatial (image) and sequential (text) tasks.
The low effective rank on images (4 vs 51 for text) supports the rank diagnostic conjecture: different domains saturate at different $r^*$, and this quantity characterises the relational complexity of the domain.

\subsection{Exp M-MAESTRO: Musical Narrative Structure}

\paragraph{Goal.}
Test whether SharedM captures musical coupling---temporal relationships between musical events in real piano performances---and whether adding text to a music model follows the sub-linear rank growth pattern.

\paragraph{Task.}
MAESTRO v3 MIDI piano performances: 89K training sequences, 64 timesteps $\times$ 6 continuous features (num\_active\_notes, mean\_pitch, pitch\_spread, mean\_velocity, bass\_note, treble\_note) at 10\,Hz.
Each token is a meaningful musical event---not a sub-phoneme spectral slice.
This event-level tokenisation gives $\M$ clean sequential structure: ``next'' means ``next musical event,'' and coupling reflects phrase structure, cadential patterns, and dynamic arcs.

\paragraph{Setup.}
$d{=}128$, $L{=}6$, $h{=}4$, 4000 steps on MPS.
Four configurations: transformer and SharedM, each music-only and music+text (Shakespeare, alternating batches).

\paragraph{Results.}

\begin{table}[h]
\centering\small
\begin{tabular}{@{}lrrrrc@{}}
\toprule
\textbf{Config} & \textbf{Params} & \textbf{Music MSE} & \textbf{vs TF} & \textbf{Text Loss} & \textbf{M Rank} \\
\midrule
TF music-only       & 1.20M & 0.01628 & ---     & ---   & --- \\
SharedM music-only  & 1.03M & 0.01568 & $-3.7\%$ & ---   & 2 \\
TF music+text       & 1.27M & 0.01644 & $+1.0\%$ & 0.500 & --- \\
SharedM music+text  & 1.11M & 0.01675 & $+2.9\%$ & 0.497 & 5 \\
\bottomrule
\end{tabular}
\caption{MAESTRO piano prediction results ($d{=}128$, 4000 steps).  SharedM beats the transformer on music-only by $3.7\%$ with $14\%$ fewer parameters.  Music coupling requires only rank 2---simpler than images (4) or text at larger scale.  Adding text grows rank from 2 to 5: sub-linear, with 3 new dimensions accommodating text coupling.}
\label{tab:maestro}
\end{table}

\paragraph{Key findings.}
\begin{enumerate}[nosep]
    \item \textbf{SharedM beats transformer on music by $3.7\%$} with $14\%$ fewer parameters.  Musical coupling (phrase boundaries, dynamic arcs, cadential patterns) is well-captured by a single shared $\M$, consistent with the image and text results.

    \item \textbf{Music coupling needs only rank 2.}  This is the lowest rank of any modality tested---simpler than images (4) or audio spectrograms (6).  Event-level tokenisation at 10\,Hz produces clean temporal structure where the dominant coupling pattern is ``attend to recent musical context,'' a rank-1 pattern, plus a secondary pattern for longer-range phrase structure.

    \item \textbf{Adding text grows rank from 2 to 5.}  The increase is sub-linear: 3 new dimensions for text, not the full rank-2 + rank-$k$ sum.  This is consistent with the three-modality rank compression in Table~\ref{tab:confirmed-rank}.

    \item \textbf{Text performance nearly identical.}  Both architectures achieve ${\sim}0.50$ text loss (0.500 TF vs 0.497 SharedM), confirming that sharing $\M$ across music and text does not degrade either modality.
\end{enumerate}

\paragraph{Interpretation.}
MAESTRO extends the SharedM advantage to a fourth modality type (structured music events, vs images, spectrograms, and text).  The rank-2 coupling geometry confirms that event-level tokenisation---where each token carries musical meaning---gives $\M$ simpler structure to discover than raw spectral frames.  The sub-linear rank growth from music-only to music+text (2$\to$5) adds further evidence that $\M$ finds shared coupling dimensions across modalities.

\subsection{Exp M3: Rank Saturation Across Modalities}

\paragraph{Goal.}
Measure $r^*$ for text-only vs text+image.

\paragraph{Setup.}
Train SharedM with $A, B \in \R^{d \times r}$ for $r \in \{4, 8, 16, 32, 64\}$ on:
\begin{enumerate}[nosep]
    \item Text-only (TinyStories captions)
    \item Image+text (CIFAR-10 patches + captions)
\end{enumerate}
Plot loss vs $r$ for both conditions.
The gap $r^*_{\text{multimodal}} - r^*_{\text{text}}$ is the measurement.

\subsection{Exp M4: Audio Next-Frame Prediction}

\paragraph{Goal.}
Test whether $\M$ learns temporal coupling for audio---the audio analogue of next-patch prediction.

\paragraph{Setup.}
Speech Commands dataset (10K train, 2K test).
Audio clips re-patched into 64 frames $\times$ 50 dimensions.
Raster (temporal) order with causal mask: predict next frame from previous frames.
Same architecture as M2, shared $\M$ handles temporal instead of spatial coupling.

\paragraph{Prediction.}
SharedM works for audio with the same architecture, no modality-specific changes.
The coupling rule ``recent frames relate to current frame'' is structurally similar to ``adjacent patches relate to current patch.''

\subsection{Exp M5: Cross-Modal --- One M for Images AND Audio (Ready)}

\paragraph{Goal.}
The strongest test: train a single model on both image patches and audio frames simultaneously.
Same $\M$, same response heads.
Alternate image and audio batches during training.

\paragraph{Setup.}
CIFAR-10 (64 patches $\times$ 48-dim) and Speech Commands (64 frames $\times$ 50-dim).
Per-modality input/output projections map to/from a shared $d$-dimensional space.
The coupling backbone ($\M$ or per-layer Q,K) is shared across modalities.

\paragraph{Models.}
Same three: Transformer, SharedM, ORN.
Evaluate separately on image MSE and audio MSE.

\paragraph{Prediction.}
SharedM handles both modalities with one coupling geometry.
Combined MSE within $5\%$ of the transformer.
If confirmed, $\M$ is modality-general---a universal relational geometry.

\paragraph{Script.} \texttt{exp\_m\_crossmodal.py} --- written and dry-run tested.

\subsection{Future Experiments}

\begin{description}[nosep, leftmargin=1.5em]
    \item[M6: Audio next-frame (standalone)] Same as M2 but audio-only. Baseline for M5 comparison.
    \item[M7: Transfer --- freeze M from images, fine-tune on audio.] The multimodal LP1. If frozen image-trained $\M$ transfers to audio, the geometry is modality-general.
    \item[M8: Semigroup analysis on cross-modal M.] Compare spectral structure of $\M$ trained on images+audio vs text-only vs image-only. Do multimodal models develop richer spectra?
    \item[M9: Image captioning.] Patches $\to$ text tokens autoregressively. Cross-modal coupling in one sequence.
    \item[M10: Video prediction.] Frame sequences at longer horizons. Moving MNIST or similar.
    \item[M11: Speech $\to$ transcription.] Cross-modal ASR-like task with paired audio-text.
    \item[M12: Three-modal simultaneous.] Image + audio + text, one $\M$. How does $r^*$ scale with number of modalities?
    \item[M13: Cross-modal transfer at 50M.] Freeze $\M$ trained on OWT text, fine-tune on images with frozen coupling. The flagship multimodal transfer result.
    \item[M14: Flexible masking and generation.] See Section~\ref{sec:generation} below.
\end{description}

\subsection{From Completion to Generation}\label{sec:generation}

The progressive completion figure (Figure~\ref{fig:m2-progressive}) is already suggestive: from 9 patches (the top-left corner), the model builds a recognisable dog.
By 33 patches (half the image), the composition is correct.
The model is \emph{already generating}---it is generating from partial context rather than from noise.

The step from completion to generation is a change of masking, not architecture.
The coupling geometry $\M$ does not care about raster order; it couples tokens based on their position in the field.
Change the attention mask, change the generation task:

\begin{description}[nosep, leftmargin=1.5em]
    \item[Sparse completion:] Give 8 random patches out of 64, predict the other 56. Tests whether $\M$ can reconstruct from scattered spatial context---closer to how vision works (sparse fixations, not raster scan).
    \item[Outpainting:] Give the centre, predict the edges. Tests extrapolation from dense context.
    \item[Inpainting:] Give the edges, predict the centre. Tests interpolation from surrounding context.
    \item[Unconditional generation:] Give only a learned seed token, predict all 64 patches. The orbital dynamics must generate from nothing---transporting a compact seed through $\M$'s geometry to produce a coherent image. This is the ORN as a generative model.
\end{description}

All four tasks use the same architecture. Only the mask changes.

\paragraph{Adaptive generation order.}
Raster scan is arbitrary---the eye does not scan images left-to-right.
Biological saccades jump to \emph{informative} regions first.
An ORN with flexible masking could learn not just what to predict but \emph{where to look next}: predict the easy patches first (uniform sky---high confidence, small correction), then the hard patches (object boundaries---low confidence, large correction).
The perturbative gate provides a natural confidence signal: patches where the gate is low can be predicted with the mean field alone; patches where the gate opens need more context before prediction.

This connects the ORN to active inference~\cite{friston2010free}: the agent selects observations (patch predictions) to maximise information gain, guided by the coupling geometry $\M$ (the generative model) and the perturbative gate (the precision/confidence signal).

\paragraph{Comparison with diffusion.}
A diffusion model generates CIFAR-10 images in ${\sim}50$--$1{,}000$ denoising steps, each a full network pass.
An ORN generates in a single forward pass of $L$ layers ($L{=}8$ in M2).
If the perturbative corrections can learn to recover the instance-specific detail that the mean field ($\M$) averages over, the ORN produces diffusion-quality results at transformer-speed inference---no iterative denoising loop.

This is not yet demonstrated.
The M2 reconstructions are blurry (MSE loss, 32$\times$32 resolution, corrections barely active).
But the mechanism is sound: $\M$ handles structure, corrections handle detail, the gate schedules the refinement.
A concrete test (M14): unconditional CIFAR-10 generation from learned seeds, with the gate amplitude and correction quality measured against resolution.
Competitive FID ($< 50$) would confirm the mechanism works; competitive with diffusion ($< 5$) would require higher resolution, richer corrections, and likely a multi-scale approach (VAR-style coarse-to-fine within the orbital framework).

% ======================================================================
\section{From Mean Field to Sharp Predictions}

\subsection{Why M Produces Blurry Images}

The M2 reconstructions (Figure~\ref{fig:m2-progressive}) capture structure---object shape, colour, composition---but the predictions are soft rather than sharp.
This is not a bug of the architecture; it is a \emph{consequence of what $\M$ is}.

$\M$ encodes the \textbf{mean-field coupling geometry}: the average relational structure across all images.
``Patches to the right of sky-coloured patches are probably also sky-coloured'' is a coupling rule that holds on average.
When $\M$ generates a prediction, it produces the expectation over all images consistent with the context so far---a blurry composite of all possible continuations.
This is the same phenomenon that makes MSE-trained autoencoders produce blurry reconstructions: optimising the average prediction hedges across modes rather than committing to one.

In physics: $\M$ gives the \textbf{classical path}---the saddle-point solution of the action.
The classical path captures the dominant trajectory but not the fluctuations.
A classical field theory of electromagnetism tells you the average field, not the specific photon.

A crucial distinction (see companion document \textit{Path Integrals and the ORN}): the ORN does not compute the path integral $T^L$ (which would average over all futures, converging to the stationary distribution).
Instead, the deeper layers \emph{refine} the context-specific prediction---the residual connection preserves the starting state, and each layer sharpens rather than averages.
The semigroup provides the geometric scaffold (composable coupling kernels, CK ${\sim}0.99$); the response heads and corrections provide the context-specific selection.
The path integral computes what \emph{can} happen; the ORN computes what \emph{will} happen given this context.

\subsection{Perturbative Corrections as Detail Recovery}

The ORN architecture already contains the mechanism for sharpening: the \textbf{perturbative corrections}.
\begin{equation}
    \x = \x_{\text{orbital}} + g(\x) \cdot \delta(\x)
\end{equation}

In the image prediction context:
\begin{description}[nosep, leftmargin=1.5em]
    \item[$\x_{\text{orbital}}$:] The mean-field prediction. ``Something dog-coloured goes here, given the spatial context.'' Correct in structure, blurry in detail.
    \item[$\delta(\x)$:] The instance-specific correction. ``This particular dog has a white chest and brown ears, given the specific patches already predicted.'' High-frequency detail that the mean field averages over.
    \item[$g(\x)$:] The gate controlling correction amplitude. When context is weak (few patches seen), $g \approx 0$ and you get the safe mean-field prediction. When context is rich (many patches seen, clear object established), $g$ opens and the correction sharpens the output.
\end{description}

This is structurally analogous to two established approaches:

\paragraph{Diffusion models.}
In a diffusion model, a denoising network iteratively removes Gaussian noise to recover a clean image.
The noise level controls how much detail the network must supply at each step: high noise $\to$ recover coarse structure; low noise $\to$ recover fine detail.
In the ORN, the orbital prediction is the ``noisy'' (blurry) estimate, and the perturbative correction is the denoising step.
The gate $g$ plays the role of the noise level---controlling how much correction (detail) to apply at each point.
The key difference: in diffusion, the noise schedule is external.
In the ORN, the gate is \emph{learned}---it reads the residual stream and decides autonomously how much correction the current context warrants.

\paragraph{Multi-scale prediction (VAR).}
Tian et al.~\cite{tian2024var} showed that predicting images coarse-to-fine (next-scale rather than next-token) dramatically improves quality (FID 18.65 $\to$ 1.73).
In the ORN, this hierarchy emerges naturally from the depth structure: early layers produce coarse coupling (broad spatial relationships), late layers refine to fine coupling (local detail).
The perturbative corrections amplify this: they are small at early layers (mean field suffices for coarse structure) and grow at late layers (detail requires instance-specific correction).

\subsection{Why the Gate Barely Opened on Images}

In M2, the perturbative gate opened to only 0.23 (vs 0.36 on text in Exp~8).
This makes sense: at 32$\times$32 resolution with 4$\times$4 patches, the spatial detail is coarse---there is not much room for instance-specific corrections to improve on the mean-field prediction.
Every dog is roughly the same colour blob at this resolution.

The prediction: \textbf{at higher resolution, the perturbative corrections should become increasingly valuable.}
Fine-grained detail (fur texture, facial features, text in images) cannot be predicted by the mean field---it requires context-specific correction.
The gate amplitude should grow with image resolution, and the MSE gap between ORN and bare SharedM should widen as resolution increases.

This is directly testable: run M2 at 64$\times$64 or 128$\times$128 (with proportionally more patches) and measure whether the ORN's corrections close the gap with the transformer more than at 32$\times$32.

\subsection{The Broader Story}

The connection between mean-field blur and perturbative sharpening is not specific to images.
In text, the mean-field prediction from $\M$ is the structurally correct continuation (``a noun probably goes here after this verb''), and the perturbative correction sharpens it to the specific word (``Lily,'' not just any noun).
The 0.449-nat ablation gap on text (Exp~8) is the cost of losing this sharpening.

The architecture proposes a division of labour that mirrors physics:
\begin{enumerate}[nosep]
    \item \textbf{$\M$ captures the symmetries}---the coupling rules that hold on average across all instances.
    \item \textbf{Response heads extract content}---what flows through the coupling pattern at each depth.
    \item \textbf{Perturbative corrections break the symmetries}---adding instance-specific detail that the mean field cannot provide, gated by context confidence.
\end{enumerate}

Images make this decomposition visible because the mean field is literally blurry and the corrections are literally sharpening.
In text the same decomposition operates but is hidden behind discrete tokens.

\subsection{The Diffusion Connection: Learned Schedules and Running Couplings}

Diffusion models and the ORN both iterate a shared operator to transport one distribution toward another: diffusion moves noise $\to$ data; the ORN moves embeddings $\to$ predictions.
But the analogy is imprecise, and the imprecision is instructive.

\paragraph{What diffusion does that ORN does not (at first glance).}
In a diffusion model, the denoising network is \emph{conditioned on the noise level $t$}.
It knows where it is in the schedule and adapts its behaviour: early steps recover coarse structure; late steps add fine detail.
The schedule is explicit.

In the ORN, $\M$ has no depth signal.
It reads the field state and produces coupling---the same $\M$ at every layer.
There is no explicit $t$.

\paragraph{But depth IS encoded---in the response heads.}
The response heads ($W_V^{(\ell)}$, $W_O^{(\ell)}$, $\text{FFN}^{(\ell)}$) are the one component that varies with depth.
They are per-layer parameters that transform the field state differently at each depth.
The residual stream at layer $\ell$ carries the cumulative signature of response heads $1$ through $\ell{-}1$: not just ``what tokens are present'' but ``what tokens are present \emph{as processed by this specific sequence of transformations}.''

When the perturbative gate reads the field state, it is reading this depth-encoded signature.
The field state after 2 layers of response-head processing looks structurally different from the field after 7 layers---not just because more context has been integrated, but because the per-layer response heads have left their fingerprint on the representation.

This means the ORN has an \textbf{implicit depth schedule}: the gate learns to vary its correction amplitude not from an explicit layer index but from the response heads' cumulative effect on the field.
The gate-entropy anticorrelation ($r = -0.95$) is evidence that this works---the gate tracks the field state's maturity, which is a direct consequence of how many response-head transformations have been applied.

\paragraph{Implicit vs explicit: which is better?}
Diffusion uses explicit noise-level conditioning: the network receives $t$ as input.
The ORN uses implicit depth conditioning: the gate reads the field state shaped by per-layer response heads.

The implicit approach has a principled advantage: the correction amplitude depends on the \emph{actual state of the representation}, not on an arbitrary layer number.
If a particular token converges quickly (function word, predictable context), the field state at layer~4 may already look ``late-stage,'' and the gate should respond accordingly.
Explicit layer conditioning would apply the same correction to all tokens at a given depth, regardless of how far each token's representation has progressed.

The implicit approach also transfers to variable-depth inference (the second paper's proposal): if the model runs for $T{=}12$ instead of $T{=}9$, the gate's behaviour at the extra layers is determined by the field state, not by a layer index that exceeds the training range.
Explicit conditioning would need extrapolation; implicit conditioning adapts naturally.

The 0.449-nat ablation gap from Exp~8 confirms the implicit schedule is working: the corrections contribute substantially, the gate opens gradually over training, and the effective contribution ($g \times \delta/h \approx 0.20$) is genuinely perturbative.
The per-layer response heads provide enough depth information through their cumulative field-state fingerprint for the gate to learn when and how much to correct.

\paragraph{The physics: running couplings, naturally.}
In quantum field theory, coupling constants ``run'' with energy scale---the effective interaction strength changes depending on the resolution at which you probe the system.
The fine-structure constant $\alpha$ is ${\sim}1/137$ at low energies but ${\sim}1/128$ at the $Z$ boson mass.
The \emph{theory} is the same; the effective coupling runs because the renormalisation group integrates out high-energy modes as you lower the resolution.

$\M$ is the theory.
The perturbative gate $g$ is the coupling constant.
The response heads at each layer are the renormalisation group steps that coarse-grain the field state.
Because the gate reads the coarse-grained field, it naturally produces a \textbf{running coupling}---the effective correction strength varies with depth, driven by how the response heads have transformed the representation.

This is not something we designed.
It is a consequence of having per-layer response heads (providing depth variation) and a context-dependent gate (reading the field state).
The running coupling emerges from the interaction of shared geometry and per-layer response, without explicit scheduling.

\paragraph{What ORN has that diffusion does not.}
The coupling geometry $\M$ has no analogue in standard diffusion.
Diffusion denoises each patch through the network's implicit coupling (attention within the denoiser).
The ORN explicitly separates the coupling question (who-attends-to-whom, governed by $\M$) from the refinement question (what correction to apply, governed by $g \cdot \delta$).

This means the ORN can learn \emph{different refinement strategies for different coupling contexts}.
A patch that couples strongly to its spatial neighbours (uniform sky region) needs different correction than a patch at an object boundary that couples to distant context.
$\M$ determines the coupling context; the perturbation responds to it.
Diffusion models must discover both coupling and refinement jointly.

% ======================================================================
\section{In-Flight Experiments}

\subsection{Frozen-M Theory Test (running on vast.ai)}

The critical transfer test: can a coupling geometry trained on one text corpus transfer to a different one without retraining?

\paragraph{Setup.}
V1's $\M$ (rank 20, trained on OpenWebText) is frozen.
Fresh response heads ($W_V$, $W_O$, FFN) are trained from scratch on FineWeb-Edu (higher quality educational text).
Control: everything trained from scratch on FineWeb-Edu.
Architecture: full ORN ($d{=}512$, $L{=}24$, $H{=}8$, $d_{\text{corr}}{=}64$), 91M total parameters, 524K frozen coupling parameters.
Running on vast.ai RTX 4090, ${\sim}40$h total for both runs.

\paragraph{Prediction.}
If frozen $\M$ matches from-scratch within $5\%$, the coupling geometry is generic and transferable---the strongest evidence yet that memoisation captures a universal relational rule, not a corpus-specific pattern.
If frozen $\M$ is $> 10\%$ worse, $\M$ needs co-training with the response heads and the universality claim is weakened.

\subsection{4-Modality Exploration}

We ran a 4-modality experiment (CIFAR-10 images, MAESTRO piano, ESC-50 environmental sounds, TinyStories text) at $d{=}128$, $L{=}6$, 3000 steps on MPS---the regime where $\M$ matters (46\% gap vs random frozen $\M$ at $d{=}96$).
The experiment trains 11 configurations: each modality alone (4 runs), all four simultaneously, a TF baseline on all four, and 5 modality pairs.
We first state the conjectures made before running, then report results and evaluate each conjecture.

\paragraph{Conjectures (pre-registered).}
\begin{enumerate}[nosep]
    \item \textbf{Single-modality ranks form a coupling complexity ordering.}
    From prior results: music $\approx 2$, images $\approx 4$, text $\approx 2$--$4$.
    We predicted environmental sounds rank highest ($> 5$) because 50 diverse classes span a wider range of temporal structures than one instrument or one image distribution.

    \item \textbf{All-4 rank is sub-linear.}
    If single ranks sum to ${\sim}15$, we predicted all-4 rank of 4--8 (2--4$\times$ compression).
    The 3-modality result (rank 3 vs sum 10, 3.3$\times$ compression) set the baseline.

    \item \textbf{Modality pairs reveal two kinds of coupling geometry.}
    Music + text should share the most (both are sequential narrative).
    Images + environmental sounds should share (both have spatial/spectral structure).

    \item \textbf{SharedM's advantage grows with modality count.}
    On 4 modalities, SharedM should win by more because sub-linear rank gives it a budget advantage.
\end{enumerate}

\paragraph{Single-modality ranks.}
Each modality was trained in isolation to measure $\M$'s crystallised rank.

\begin{center}\small
\begin{tabular}{@{}lrrp{6cm}@{}}
\toprule
\textbf{Modality} & \textbf{Rank} & \textbf{Loss} & \textbf{Crystallisation} \\
\midrule
Music (MAESTRO) & 2 & 0.0163 & Rank $39 \to 2$ by step 1000 (fastest) \\
Images (CIFAR-10) & 4 & 0.0507 & Rank $14 \to 4$ by step 3000 (gradual) \\
Text (TinyStories) & 4 & 0.4918 & Rank $7 \to 4$ by step 1000 \\
Env sounds (ESC-50) & 5 & 8.9955 & Rank $14 \to 5$ by step 3000 (gradual) \\
\bottomrule
\end{tabular}
\end{center}

The coupling complexity ordering is: music (2) $<$ images (4) $=$ text (4) $<$ environmental sounds (5).
Music crystallises fastest and to the lowest rank, consistent with MAESTRO's relatively homogeneous temporal structure (single instrument, regular onset patterns).
Environmental sounds, with 50 diverse classes spanning very different temporal envelopes, require the highest rank.

\paragraph{Pair ranks and quality costs.}
For each modality pair, we report the shared $\M$'s crystallised rank, the sum of single-modality ranks, the resulting compression ratio, and the per-modality quality cost relative to the single-modality baseline.

\begin{center}\small
\begin{tabular}{@{}lrrrp{5cm}@{}}
\toprule
\textbf{Pair} & \textbf{Rank} & \textbf{$\Sigma$ singles} & \textbf{Compression} & \textbf{Quality cost} \\
\midrule
mus+txt & 4 & 6 & $1.5\times$ & mus $+2.9\%$, txt $+1.7\%$ \\
img+mus & 5 & 6 & $1.2\times$ & img $+36.9\%$, mus $+6.8\%$ \\
img+env & 5 & 9 & $1.8\times$ & img $+52.3\%$, env $+63.4\%$ \\
env+txt & 3 & 9 & $3.0\times$ & env $+59.4\%$, txt $+3.2\%$ \\
mus+env & 2 & 7 & $3.5\times$ & mus $-0.7\%$, env $+47.9\%$ \\
\bottomrule
\end{tabular}
\end{center}

The critical insight: \textbf{rank compression $\neq$ quality preservation.}
The mus+env pair achieves rank 2 (best compression, $3.5\times$), but environmental sounds degrade $47.9\%$.
The mus+txt pair achieves rank 4 (moderate compression, $1.5\times$) with under $3\%$ degradation on both modalities.
The right metric is rank compression \emph{at acceptable quality loss}, not rank compression alone.

Music and text are the most \emph{compatible} pair: rank $4 = \max(2, 4)$, barely any quality cost.
Music's rank-2 coupling fits inside text's rank-4 subspace without competition.
More generally, music is robust to sharing---it never degrades more than $6.8\%$ in any pairing.
Its coupling is so compact that it does not compete for $\M$'s capacity.

Images and environmental sounds are sensitive: they degrade significantly whenever sharing $\M$.
This is precisely where perturbative corrections should help---recovering per-modality detail that $\M$'s compression loses.
The corrections are context-gated and the modality token is in the residual stream, so the gate can learn ``activate for environmental sounds, stay quiet for music.''
This predicts modality-dependent gate magnitudes, with env sounds and images showing the largest gate values.

\paragraph{All-4 results.}
Training all four modalities simultaneously, $\M$ crystallises to rank 4 (sum of singles $= 15$, compression $3.8\times$).
Per-modality quality costs relative to single-modality baselines: images $+63.0\%$, music $+6.3\%$, env sounds $+112.4\%$, text $+7.0\%$.
The pattern is consistent: music and text tolerate sharing; images and env sounds suffer.

\paragraph{Three architectures on all-4.}
We compare three architectures at matched parameter count ($d{=}128$, $L{=}6$, 3000 steps on MPS), all trained on the same interleaved 4-modality data: a standard transformer baseline, bare SharedM, and the full ORN (SharedM + perturbative corrections).

\begin{center}\small
\begin{tabular}{@{}lrrrrrr@{}}
\toprule
& \textbf{Images} & \textbf{Music} & \textbf{Env sounds} & \textbf{Text} & \textbf{M rank} & \textbf{Gates} \\
\midrule
Transformer & 0.0801 & 0.0174 & 17.84 & 0.5215 & --- & --- \\
SharedM     & 0.0828 & 0.0174 & 19.41 & 0.5258 & 4 & --- \\
ORN         & 0.0846 & 0.0167 & 17.30 & 0.5328 & 4 & \scriptsize img\,0.09 mus\,0.05 env\,0.21 txt\,0.21 \\
\bottomrule
\end{tabular}
\end{center}

The full ORN \emph{beats} the transformer on music ($-3.9\%$) and environmental sounds ($-3.1\%$), while losing on images ($+5.7\%$) and text ($+2.2\%$).
Bare SharedM is within $9\%$ of the transformer on every modality and matches it on music ($+0.2\%$).

\paragraph{Quality recovery analysis.}
To quantify how effectively corrections compensate for $\M$'s compression loss, we measure \emph{recovery}: how much of SharedM's gap to the transformer is closed (or exceeded) by corrections.

\begin{center}\small
\begin{tabular}{@{}lrrr@{}}
\toprule
\textbf{Modality} & \textbf{SM gap vs TF} & \textbf{ORN gap vs TF} & \textbf{Recovery} \\
\midrule
Images        & $+3.4\%$  & $+5.7\%$  & $-65\%$ (worse) \\
Music         & $+0.2\%$  & $-3.9\%$  & $2623\%$ (beats TF) \\
Env sounds    & $+8.8\%$  & $-3.1\%$  & $135\%$ (beats TF) \\
Text          & $+0.8\%$  & $+2.2\%$  & $-163\%$ (worse) \\
\bottomrule
\end{tabular}
\end{center}

Corrections selectively compensate where $\M$'s compression hurts most.
Environmental sounds, which suffered $+8.8\%$ under bare SharedM, are brought \emph{below} the transformer baseline by $3.1\%$.
Music, already near-parity, is pushed to a clear win.
Images and text, however, get \emph{worse} with corrections at this scale---the limited correction capacity (a single small MLP per layer) is allocated to the modalities where it helps most, leaving images and text underserved.

\paragraph{Gate magnitudes are modality-adaptive.}
The gates reveal the network's correction allocation strategy.
Comparing single-modality training to the all-4 setting:

\begin{center}\small
\begin{tabular}{@{}lrrr@{}}
\toprule
\textbf{Modality} & \textbf{Single} & \textbf{All-4} & \textbf{Change} \\
\midrule
Images  & 0.287 & 0.093 & $-0.194$ (closes) \\
Music   & 0.032 & 0.052 & $+0.020$ (barely) \\
Env     & 0.101 & 0.209 & $+0.108$ (opens) \\
Text    & 0.116 & 0.205 & $+0.089$ (opens) \\
\bottomrule
\end{tabular}
\end{center}

Image gates \emph{close} in the multi-modal setting ($-0.194$), while env and text gates open substantially.
This is rational triage: image corrections were large in isolation (gate 0.287) but are deprioritised when the network must share correction capacity across modalities.
Environmental sounds, which had the largest SharedM gap ($+8.8\%$), get the widest gate opening ($+0.108$), and this investment pays off (recovery to $-3.1\%$ vs transformer).

\paragraph{M rank is invariant to corrections.}
$\M$ crystallises to rank~4 with or without perturbative corrections.
Corrections do not alter the coupling geometry---they compensate in the value stream only.
This confirms the architectural decomposition: $\M$ governs \emph{who couples to whom}; corrections govern \emph{what to do about it} on a per-instance basis.

\paragraph{Multi-task degradation is a training budget problem, not a correction problem.}
The quality cost of training on all four modalities simultaneously (each modality sees ${\sim}25\%$ of training steps) is nearly identical for SharedM and ORN: ${\sim}63$--$70\%$ on images, ${\sim}7\%$ on text, ${\sim}112$--$117\%$ on env sounds, ${\sim}3$--$6\%$ on music.
Corrections do not fix training budget starvation---they fix $\M$'s coverage gap within the budget each modality receives.

\paragraph{Conjecture evaluation.}
\begin{itemize}[nosep]
    \item \textbf{Conjecture 1} (coupling complexity ordering): \textbf{Confirmed.}
    The ordering env sounds (5) $>$ images (4) $=$ text (4) $>$ music (2) matches the prediction that diverse environmental sounds require the highest coupling rank.

    \item \textbf{Conjecture 2} (sub-linear rank): \textbf{Confirmed.}
    All-4 rank $= 4$ vs sum of singles $= 15$, giving $3.8\times$ compression.
    This exceeds the 3-modality baseline ($3.3\times$) and falls squarely in the predicted 2--4$\times$ range.

    \item \textbf{Conjecture 3} (mus+txt most shared): \textbf{Not confirmed as stated.}
    The mus+env pair achieves the greatest rank compression ($3.5\times$, rank 2), not mus+txt.
    However, mus+txt is the best \emph{compatible} pair---highest compression with acceptable quality ($< 3\%$ cost on both modalities).
    The conjecture was wrong about which pair compresses most but right about which pair works best in practice.

    \item \textbf{Conjecture 4} (SharedM advantage grows with modality count): \textbf{Partially confirmed.}
    The full ORN beats the transformer on 2 of 4 modalities (music and env sounds), demonstrating that corrections can convert $\M$'s compression into a net advantage.
    On images and text, corrections at this scale are insufficient.
    The prediction is that at larger $d$ (more correction capacity) and longer training (resolving budget starvation), the advantage extends to all modalities.
\end{itemize}

\paragraph{Scaling outlook.}
At $d{=}128$ with 3000 steps, correction capacity is stretched across four modalities.
At V2 scale ($d{=}512{+}$, longer training), we expect three improvements.
First, correction capacity scales with $d$---the correction MLP at $d{=}512$ has $16\times$ the capacity of $d{=}128$, sufficient to serve all modalities simultaneously.
Second, longer training resolves budget starvation: each modality sees enough data to converge, eliminating the ${\sim}63$--$112\%$ multi-task degradation.
Third, image representations may need rethinking---we invested research in audio tokenisation and it paid off (env sounds benefit most from corrections), but images still use naive raster-scan patches.
Saccade-ordered or multi-resolution patch strategies could improve image coupling compatibility with the shared $\M$.
The goal: at scale, match or beat the transformer on \emph{all} modalities with fewer coupling parameters, creating a world model where adding modalities strengthens rather than dilutes each one.

\subsection{Crossover Trajectory Confirmed}

To test whether the ORN--transformer gap narrows with training, we ran both architectures on all four modalities for 8000 steps (2000 per modality), taking snapshots at 3000, 5000, and 8000 steps.

\begin{center}
\begin{tabular}{lrrrr}
\toprule
\textbf{Steps} & \textbf{Images} & \textbf{Music} & \textbf{Env} & \textbf{Text} \\
\midrule
3000 & $+2.4\%$ & $+4.0\%$ & $+11.8\%$ & $+0.4\%$ \\
5000 & $+1.2\%$ & $-1.5\%$ & $-17.7\%$ & $-1.3\%$ \\
8000 & $-2.6\%$ & $+9.1\%$ & $-4.2\%$ & $+2.4\%$ \\
\bottomrule
\end{tabular}
\end{center}

\noindent
Positive values mean ORN trails the transformer; negative values mean ORN leads.

\paragraph{Images cross over.}
ORN trails by $+2.4\%$ at 3000 steps, narrows to $+1.2\%$ at 5000, and overtakes at $-2.6\%$ by 8000.
This confirms the ``world modeller starts slow'' hypothesis: SharedM plus corrections needs more training to crystallise the coupling geometry, but once crystallised, it uses parameters more efficiently than per-layer memorisation.

\paragraph{Environmental sounds show the most dramatic swing.}
From $+11.8\%$ behind at 3K to $-17.7\%$ ahead at 5K, settling at $-4.2\%$ ahead at 8K.
The large oscillation suggests the coupling geometry undergoes a phase transition around 4--5K steps where $\M$ locks in the environmental sound structure.

\paragraph{Music and text oscillate without clear trend.}
These modalities are already well-served by $\M$'s compact rank-2 coupling, so neither architecture has a systematic advantage---both are near-optimal for the available budget.

\paragraph{Interpretation.}
The crossover confirms the central prediction of the ORN multimodal theory.
A transformer memorises per-layer coupling rules quickly (it wins early).
The ORN must first crystallise a shared coupling geometry, \emph{then} train response heads and corrections to use it---this takes longer, but once done, the shared geometry provides a more parameter-efficient representation.

This is analogous to the single-modality V1 result (val~3.35 matching GPT-2 Small with $0.6\%$ coupling params) but now demonstrated in the multimodal setting where the advantage is amplified: $\M$'s sub-linear rank growth (rank~4 for 4 modalities vs sum~15) means each new modality is nearly free in coupling budget.

The practical implication: multimodal ORN experiments at fewer than ${\sim}2000$ steps per modality are in the pre-crossover regime and will underestimate ORN's true performance.
The V2 training plan (50B tokens) will place us well beyond the crossover point.

\subsection{Scaled 4-Modality Results ($d{=}256$, 20K Steps)}

We scaled the 4-modality experiment to $d{=}256$ with 20K training steps, comparing Transformer, ORN, and CORN at matched parameter count.

\begin{center}\small
\begin{tabular}{@{}lrrrrc@{}}
\toprule
& \textbf{Images} & \textbf{Music} & \textbf{Env sounds} & \textbf{Text} & \textbf{M rank} \\
\midrule
Transformer & 0.0454 & 0.0173 & 6.39 & 0.4852 & --- \\
ORN         & 0.0464 & 0.0153 & 5.58 & 0.5079 & 17 \\
CORN        & 0.0515 & 0.0149 & 5.42 & 0.4807 & 11 \\
\bottomrule
\end{tabular}
\end{center}

CORN wins on three of four modalities: music ($-13.8\%$), environmental sounds ($-15.3\%$), and text ($-0.9\%$).
Only images trail ($+13.4\%$).
This confirms the crossover trajectory predicted at $d{=}128$: with sufficient scale and training, the shared coupling geometry outperforms per-layer memorisation.

\paragraph{Rank compression at scale.}
CORN's $\M$ compresses to rank~11, down from ORN's rank~17.
The semigroup consistency loss concentrates coupling energy into fewer modes, achieving tighter spectral structure without sacrificing quality---in fact, improving it on 3/4 modalities.

\paragraph{Semigroup generation: $t = \sigma$ is optimal.}
Parallel experiments on diffusion-style generation via CORN's $t$-scaled semigroup dynamics confirmed that $t = \sigma$ (coupling strength equals noise level) is the optimal mapping between semigroup time and diffusion noise.
This was verified against 5 alternative mappings ($t = \sigma^2$, $t = \log\sigma$, linear schedules, cosine schedules).
The result is natural: the semigroup's eigenspectrum \emph{is} the noise schedule, and the identity mapping $t = \sigma$ respects this correspondence directly.

% ======================================================================
\subsection{Exp M-V3: Large-Scale Text-to-Multimodal Transfer (605M, April 2026)}\label{sec:v3-transfer}

The largest-scale frozen-M transfer test to date: a 605M ORN trained exclusively on text is adapted to image captioning and audio captioning with $0.83\%$ of parameters trainable.

\paragraph{Setup.}
A 605M ORN ($d{=}2048$, $L{=}32$, GQA 32/8, 8$\times$ wide shared FFN, $d_{\text{ff}}{=}44{,}032$) was trained on FineWeb-Edu for 819M tokens.
The backbone was then frozen---SharedM ($A$, $B$) and the wide shared FFN---and lightweight adapters were trained on two simultaneous captioning tasks:
\begin{itemize}[nosep]
    \item \textbf{Image captioning:} COCO train2017 (118K images), 2D patch encoder, ViT-style 16$\times$16 patches, learned 2D position embeddings
    \item \textbf{Audio captioning:} AudioCaps + Clotho, 20 spectrogram patches $\times$ 1600-dim, temporal position embeddings
\end{itemize}
Trainable components: LoRA rank~16 on $W_V$, $W_O$ (per-layer); patch encoder; audio encoder; caption token embeddings.
Total trainable: 4,988,928 parameters ($0.83\%$ of 600M).
Training: 50K steps, cosine LR with warmup, effective batch size 4 (grad accum 4), alternating 75/25 image/audio.
Hardware: Vast.ai RTX 4090, 12{,}435s (${\sim}3.5$h).

\paragraph{Results.}

\begin{center}\small
\begin{tabular}{@{}lrr@{}}
\toprule
\textbf{Metric} & \textbf{Image} & \textbf{Audio} \\
\midrule
Val loss (final, step 50K) & 3.190 & 3.085 \\
Val loss (step 47.5K) & 3.333 & 2.968 \\
\bottomrule
\end{tabular}
\end{center}

$\M$'s spectral structure remained stable throughout: top-SV~$= 73.73$, rank-90\%~$= 73$ (frozen by design, confirmed unchanged).

\paragraph{Generation samples (step 50K).}

\begin{center}\small
\begin{tabular}{@{}lp{9cm}@{}}
\toprule
\textbf{Source} & \textbf{Text} \\
\midrule
IMG ground truth & \textit{A cat laying on top of a wooden computer desk.} \\
IMG generated & A photo of a man holding a frisbee in his hand. \\
\midrule
AUD ground truth & \textit{People are talking as gun is fired} \\
AUD generated & A man speaks and a dog barks and a woman speaks in the background. \\
\bottomrule
\end{tabular}
\end{center}

\paragraph{What this confirms.}
\begin{enumerate}[nosep]
    \item \textbf{Frozen transfer works.}  A text-trained backbone adapts to two new modalities simultaneously with $0.83\%$ parameter overhead.  Training loss descended from ${\sim}5.5$ (image) and ${\sim}6.2$ (audio) at step~1 to final values around 3 for both modalities.
    \item \textbf{Conditioning is real.}  Real images produce scene-plausible captions; blank (zero) images produce incoherent outputs.  The patch encoder is genuinely injecting modality-specific signal into the residual field.
    \item \textbf{Audio adapts faster than image.}  Audio loss dropped to 2.35 by step~400 (still in warmup); image loss was slower.  Consistent with audio captions being simpler (``the sound of X'') and requiring less visual grounding.
    \item \textbf{$\M$ rank is stable under multimodal fine-tuning.}  Rank-90\%~$= 73$ throughout.  Freezing $\M$ does not destabilise training.
\end{enumerate}

\paragraph{What this does not yet confirm.}
\begin{enumerate}[nosep]
    \item \textbf{Image-specific generation.}  The model produces scene-plausible captions (grammatical, contextually reasonable) but not image-specific ones---``man holding a frisbee'' for a cat on a desk.  The residual field carries a signal that ``an image is present'' but does not fully bind image content to caption tokens.
    \item \textbf{$\M$ as cross-modal coupling authority.}  The current result is consistent with a weaker explanation: the patch encoder provides a genre signal (``this is a captioning task'') and the frozen $\M$+FFN generates appropriate language, with LoRA doing the content binding.  Whether $\M$ mediates the cross-modal coupling or merely provides the language generation backbone is not yet distinguished.
\end{enumerate}

\paragraph{Interpretation and next steps.}
The result is a proof of concept for the lightweight adaptation claim---not for the coupling authority claim.
The image-specific generation gap is precisely what the MoE design addresses: with dedicated modality experts routed by $\M$'s geometry, image content would activate a visual expert rather than relying on the frozen language FFN.

Two diagnostic experiments can close the gap without additional GPU training:
\begin{itemize}[nosep]
    \item \textbf{Cross-modal retrieval:}  Embed 100 image/audio pairs through the frozen backbone; measure whether cosine similarity clusters semantically across modalities.  If it does, $\M$ is the coupling authority.
    \item \textbf{Rank analysis:}  Measure rank-90\% of $\M$'s residual representations on text vs image vs audio inputs from the fine-tuned checkpoint.  Sub-linear growth would confirm the shared geometry claim at 605M scale.
\end{itemize}

A stronger GPU experiment---unfreezing $\M$ during multimodal training and measuring whether image-specific generation improves---would directly test whether $\M$ is the bottleneck.
If trainable $\M$ produces image-specific captions and frozen $\M$ does not, that is the most informative result possible: $\M$ is the coupling authority, and the right fix is to let it adapt to the new modalities.

% ======================================================================
\section{Open Questions Requiring Experiments}

\begin{description}[nosep, leftmargin=1.5em]
    \item[Exp M5: Does one M handle images + audio? (ready)]
    The cross-modal experiment is scripted.
    If SharedM matches the transformer on both modalities simultaneously, the coupling geometry is modality-general.

    \item[Exp M7: Does frozen image-trained M transfer to audio?]
    The strongest modality-transfer test.
    If frozen $\M$ from image training adapts to audio with only response-head fine-tuning, the geometry encodes abstract relational structure (``nearby things relate'') independent of whether ``nearby'' means spatial or temporal.

    \item[Exp M5-10M: Cross-modal at larger scale]
    M5 at 2M params is a feasibility check.
    At 10M, does the cross-modal advantage grow?
    The coupling redundancy argument predicts it should.

    \item[Exp M14: Flexible masking and generation]
    Sparse completion, inpainting, outpainting, unconditional generation---all from the same architecture.
    The perturbative gate as confidence signal for adaptive generation order.

    \item[Mean-field sharpening at higher resolution]
    The perturbative corrections barely activated on 32$\times$32 CIFAR (gate 0.23).
    At 64$\times$64 or 128$\times$128, do they activate more and close the blurriness gap?
    This directly tests the mean-field $\to$ correction decomposition on images.

    \item[Path integral training on images (Exp PI-ORN)]
    Does multi-horizon prediction (depth $L$ predicts patch $L$ steps ahead) improve image completion coherence?
    See companion document \textit{Path Integrals and the ORN} for the theoretical framework.
\end{description}

\section{Modality Representation: Playing to the Geometry}

Our current experiments use the simplest modality interface: a linear projection from modality-specific dimension to $d$.
This is fair for benchmarking (identical interface for transformer and ORN) but not necessarily optimal for either.

The ORN's advantage is that $\M$ encodes relational structure.
The richer the relational signal in the input, the more $\M$ has to work with.
How we represent each modality determines how much relational structure is available for $\M$ to discover.

\paragraph{Audio patching matters.}
Our initial audio results were weak ($+9.3\%$ behind transformer) because the re-patching interleaved frequency and time dimensions, confusing the coupling geometry.
Restoring the original temporal structure (16 frames $\times$ 200-dim, clean time axis) gives $\M$ a sequence where ``next'' means ``next in time''---the coupling $\M$ should naturally discover temporal adjacency, periodicity, and hierarchical composition (frames $\to$ phonemes $\to$ words).

\paragraph{M12 three-modal results (audio fix).}
With the corrected audio representation, we ran the full three-modal experiment (M12) at 13M parameters, $d{=}192$, $L{=}8$, 6000 steps.
The results confirm that audio patching was the primary confound:

\begin{center}\small
\begin{tabular}{@{}lrrrr@{}}
\toprule
\textbf{Model} & \textbf{Image MSE} & \textbf{Audio MSE} & \textbf{Text CE} & \textbf{Combined} \\
\midrule
Transformer (13.3M) & 0.0566 & 6.6128 & 2.9629 & 9.6323 \\
ORN (13.0M) & 0.0526 & 6.7574 & 3.0296 & 9.8396 \\
Delta & $-7.0\%$ & $+2.2\%$ & $+2.3\%$ & $+2.2\%$ \\
\bottomrule
\end{tabular}
\end{center}

ORN wins on images ($-7.0\%$) with 2.7\% fewer parameters, confirming the coupling compression advantage for spatial modalities.
Audio and text are slightly behind ($+2.2\%$ each), consistent with the ORN's per-layer response heads needing more capacity for temporal and sequential modalities at this scale.
The combined loss is within $2.2\%$---well inside the $5\%$ threshold for the multimodal memoisation claim.
Perturbative gate values (image 0.042, audio 0.094, text 0.109) show the expected pattern: images need minimal correction (spatial coupling is well-captured by $\M$), while audio and text require more perturbative refinement.

\paragraph{Multi-scale input.}
Images could be represented at multiple resolutions: coarse patches for global composition, fine patches for local detail.
This maps onto the mean-field + correction decomposition: $\M$ handles coarse coupling (``sky region connects to sky region''), corrections sharpen fine detail (``edge between dog and background'').
Multi-scale input has no analogue benefit for per-layer Q,K architectures, making it a potential ORN advantage.

\paragraph{Learned modality projections.}
A small per-modality MLP mapping raw input into ``M-friendly'' space could expose relational structure more effectively than a single linear projection.
This adds parameters (to be reported) but is honest if the goal is to maximise $\M$'s coupling capacity for each modality.

\paragraph{Benchmarking principle.}
We report all results with the fair interface (same linear projection, both architectures) AND note where ORN-specific representations improve results.
The coupling compression ($8\times$ fewer coupling parameters) holds regardless of input representation.

\paragraph{Audio embedding analysis.}
To understand what structure the 16-frame $\times$ 200-dim mel spectrogram actually contains, we analysed the Speech Commands audio data before designing embedding variants.
The temporal axis decomposes into 4 distinct blocks of 4 sub-frames (intra-block correlation ${\sim}0.85$, inter-block ${\sim}0.2$), and PCA on the frequency axis shows 5 components capture 93\% of variance---the 200 frequency bins are largely redundant.
The actual information content is approximately 4 temporal events $\times$ 5 spectral features.

Six embedding variants were tested ($d{=}96$, $L{=}6$, 3000 steps, MPS).
For the directly comparable variants targeting raw 200-dim frequency:
raw 16$\times$200 baseline achieved MSE 9.81;
frequency-band projection (4$\times$50 bands projected separately, then summed) achieved MSE 9.27 ($-5.5\%$).
Block-pooled variants (averaging 4 sub-frames into 4 coarse temporal events) gave MSE 11.91 (fixed pool) and 11.53 (learnable weighted pool)---higher loss on the raw target but a cleaner task structure where each token represents a perceptual event.
PCA-compressed variants (B: PCA to 20d, MSE 91.1; D: PCA + pool, MSE 108.8) are not directly comparable due to target scaling differences.

The key finding: frequency-band projection gives a free $5.5\%$ improvement by respecting the spectral block structure.
Block pooling creates a representation where audio tokens align structurally with image patches and text tokens---all become ``sequence of events, predict next.''
The literature (AST, AudioMAE, BEATs) converges on 2D spectrogram patches ($16{\times}16$ on 128-bin mel, yielding 6--25 tokens/sec) as the standard approach, treating audio spectrograms as images---a natural fit for ORN's continuous projection.
We have now prepared richer audio data for subsequent experiments:
\begin{itemize}[nosep]
    \item \textbf{MAESTRO v3 MIDI:} 89K piano performance sequences, event-level tokenisation at 10\,Hz (64 timesteps $\times$ 6 continuous features).  Results in Section~\ref{tab:maestro}.
    \item \textbf{ESC-50:} 1{,}600 environmental sound clips as AST-style $16{\times}16$ spectrogram patches (96 patches $\times$ 256-dim).
\end{itemize}
These enable 4-modality experiments: images + music + environmental sounds + text.

\paragraph{Audio tokenisation: converging representations.}
The literature on audio transformers converges on three representation families, each with different implications for $\M$:
\begin{enumerate}[nosep]
    \item \textbf{2D spectrogram patches} (AST, AudioMAE, BEATs): $16{\times}16$ patches on 128-bin mel spectrograms, yielding 6--25 tokens/sec.  This treats audio as an image, and the standard approach for classification and understanding tasks.
    \item \textbf{Neural audio codecs} (EnCodec, SoundStream, DAC): continuous frame-level features at 50--86\,Hz.  The encoder output \emph{before} quantisation preserves the continuous geometry $\M$ needs---the codebook quantisation step discards exactly the information $\M$ operates on.
    \item \textbf{Event-level tokenisation} (MAESTRO at 10\,Hz): each token is a musically meaningful event.  Gives $\M$ the cleanest sequential structure of any audio representation.
\end{enumerate}
The semantic/acoustic token duality from AudioLM maps onto ORN's decomposition: semantic tokens capture relational structure (what $\M$ models), acoustic tokens capture surface detail (what response heads handle).  This is a natural fit for the mean-field + correction architecture.

\subsection{Design Principles for Modality Tokenisation}\label{sec:tokenisation-design}

How we tokenise each modality determines what relational structure $\M$ can discover.
A bad tokenisation hides structure; a good one exposes it.
ORN is \emph{more} sensitive to tokenisation quality than a standard transformer because it has fewer coupling parameters---a transformer with $2Ld^2$ coupling params can compensate for poor input structure by learning modality-specific rules at each layer; $\M$ with $d^2$ cannot.

\paragraph{The token rate sweet spot.}
Too fast ($>50$\,Hz): adjacent tokens are nearly identical, next-token prediction is trivial, and $\M$ wastes rank on temporal interpolation rather than structural coupling.
Too slow ($<1$\,Hz): each token is too complex, packing multiple events into one representation.
The sweet spot is 6--25 tokens/sec, where each token covers one perceptual event (a note onset, a syllable, a saccade fixation).
This is where our representations land: MAESTRO at 10\,Hz (musical events), ESC-50 at ${\sim}19$ patches/sec (spectral snapshots), images at 64 patches/image (spatial regions), text at ${\sim}5$--15 tokens/sec (words).

\paragraph{Images: saccade ordering vs raster scan.}
Raster scan (left-to-right, top-to-bottom) is arbitrary---the eye does not process images this way.
Biological saccades jump to salient regions first: face, then hands, then held object.
We tested saliency-ordered and spiral-ordered patches: both make next-patch prediction harder ($+1$--$2\%$ MSE) because the spatial continuity that makes raster prediction easy is broken.
However, saccade ordering creates a \emph{narrative} structure (``look at important thing, then next important thing'') that could align with text and music's sequential structure.
Our structured embeddings experiment showed that this alignment doesn't help for single-modality prediction, but the hypothesis remains that saccade-ordered images would share more coupling geometry with text/music than raster-ordered images in a cross-modal setting.

\paragraph{Audio: from spectral frames to perceptual events.}
Raw mel spectrogram frames at 100\,Hz are sub-phoneme slices---too fine for meaningful coupling.
Three better approaches:
\begin{itemize}[nosep]
    \item \textbf{2D spectrogram patches} (AST, AudioMAE): treat the spectrogram as an image, apply $16{\times}16$ patches.  Each patch spans ${\sim}160$\,ms $\times$ 20 mel bins---a recognisable acoustic event.  The standard approach for audio transformers.
    \item \textbf{Event-level tokenisation} (MAESTRO): each token IS a musical event (note onset, chord change, dynamic shift).  Gives $\M$ the cleanest sequential structure.  Our data: 6 continuous features per 100\,ms frame.
    \item \textbf{Neural codec features} (EnCodec, SoundStream, DAC): the encoder output \emph{before} RVQ quantisation gives continuous features at 50--86\,Hz.  The RVQ hierarchy itself mirrors ORN's mean-field/correction decomposition: first codebook captures coarse structure ($\M$'s domain), later codebooks capture acoustic detail (response heads' domain).
\end{itemize}

\paragraph{Music's multi-scale structure.}
Music has nested temporal structure spanning five orders of magnitude: notes (${\sim}50$\,ms), beats (${\sim}500$\,ms), bars (${\sim}2$\,s), phrases (${\sim}8$\,s), sections (${\sim}60$\,s).
No single token rate captures all levels.
For ORN, the hypothesis is that $\M$'s coupling rules are scale-invariant: the geometric relationship ``what follows what'' and ``what belongs together'' applies at note, phrase, and section level.
The response heads, which vary per layer, handle scale-specific content.
This is untested but suggested by the rank-2 finding for music: if $\M$ needed separate coupling rules per scale, rank would be higher.

\paragraph{Cross-modal alignment through tokenisation.}
The key insight: if all modalities are tokenised as ``sequence of perceptual events, predict next event,'' then $\M$'s coupling geometry naturally discovers whatever cross-modal structure exists.
A film score token (``tension rises'') and a text token (``the villain appeared'') are relationally similar---both set up what comes next.
Whether $\M$ discovers this depends on whether the tokenisation exposes the relational structure, not on whether we engineer an alignment loss.
CLIP and CLAP force alignment via contrastive loss.
ORN discovers alignment (or its absence) through shared coupling dynamics.

\subsection{M as a Convergence-Divergence Zone}

Damasio's convergence-divergence zone (CDZ) framework~\cite{damasio2009cdz} proposes that the brain binds multimodal representations through a hierarchy: modality-specific cortices produce sensory fragments; CDZs in association cortex store \emph{binding records} linking co-occurring features; during recall, CDZs reactivate the original fragments.

$\M$ is structurally a CDZ.
It does not represent content (that is the residual field).
It represents \emph{coupling rules}---which tokens should attend to which.
Modality-specific encoders (ViT for images, spectrogram convolutions for audio, token embeddings for text) are the sensory cortices.
The residual field is the activity pattern being coordinated.
$\M$ binds co-occurring features across modalities without storing the features themselves.

This analogy is structural, not metaphorical: a CDZ says ``these fragments co-occur''; $\M$ says ``these tokens should couple.''
Both are compact relational records, not feature stores.

\subsection{Chords, Grammar, and the Depth of Binding}\label{sec:chords}

A musical chord is a relational structure: the intervals between simultaneous notes create consonance or tension.
A chord progression is a sequence of relational structures: I--IV--V--I tells a story of departure and return.
Language has the same two levels: a grammatical chunk (``the old cat'') is a local binding of determiner, modifier, and noun; a sentence is a progression of such bindings.

The parallel is suggestive: both modalities have \emph{local coupling} (within a chunk/chord) and \emph{sequential coupling} (between chunks/chords), and both have a tension/resolution dynamic (dissonance demands resolution; a relative clause demands a main verb).
If $\M$'s geometry captures ``local binding'' and ``sequential progression'' as abstract operations, these dimensions should be shared across music and language.

However, a critical asymmetry: musical chords are \emph{explicit in the input} (multiple notes at one timestep), while grammatical chunks are \emph{emergent through depth} (individual tokens that bind across layers of processing).
A chord arrives pre-bound; a noun phrase is constructed.
This means the shared geometry, if it exists, may operate at different layers for different modalities: $\M$ applied at layer 1 to music sees pre-formed chords, while $\M$ applied at layer 3 to text sees emergent grammatical groups.
The per-layer response heads handle this naturally---the \emph{coupling rule} (``bound things attend to each other'') is the same, but the \emph{state of binding} in the residual field differs by depth and modality.

This is a reason \emph{not} to engineer representations that force alignment (e.g.\ decomposing music into individual notes to match text's sequential structure, or pre-parsing text into grammatical chunks to match music's chord structure).
If the shared geometry is real, $\M$ discovers it from natural representations at whatever depth it emerges.
If we have to engineer the representation to create overlap, we are forcing a hypothesis rather than testing it.
The 4-modality rank results (Section~\ref{sec:tokenisation-design}) test exactly this: do natural representations produce sub-linear rank growth?
They do (3.8$\times$ compression), without any representational engineering.

\subsection{Toward Orbital Modality Embeddings}

Current multimodal approaches project each modality into a shared space through alignment losses (CLIP~\cite{radford2021clip}), linear projection (LLaVA), or information bottlenecks (Q-Former).
All treat the projection and the coupling as separate mechanisms.

Orbital embeddings suggest a more principled alternative: the coupling geometry $\M$ IS the alignment mechanism.
Each modality starts from its own seed (raw patches, frames, tokens) and evolves through the same shared dynamics $\Phi$---the same operator that defines coupling.
Tokens from different modalities that are relationally similar converge to similar regions of coupling space, not because they were trained on paired data, but because $\M$'s geometry naturally places relationally similar things near each other.

This is the ORN's version of the Neural ODE~\cite{chen2018neural} connection: the orbital iteration $h_{\ell+1} = h_\ell + f(\M, h_\ell)$ is a discretised ODE with $\M$ defining the vector field.
Different modalities start from different initial conditions and evolve through the same field.
Convergence happens where the dynamics say it should.

Our embedding ablation (Section~\ref{sec:embedding-ablation}) tests this directly.  Orbital embeddings with $T{=}1$ and $T{=}3$ iterations show modest improvement ($-2.1\%$, $-3.1\%$) but are substantially outperformed by a simple modality token ($-10.0\%$).  At current scale, telling $\M$ what modality it is processing (via a learnable prefix token) is far more effective than using $\M$'s own dynamics as the embedding function.

The gate-by-modality results (image gate 0.058, audio gate 0.163, text gate 0.106) suggest that $\M$'s coverage varies by modality---a signature consistent with different modalities entering the coupling space from different directions and requiring different amounts of correction to align.  Whether orbital embeddings help at larger scale (where $\M$ has richer spectral structure) remains open.

\section{Literature: How Multimodal Models Handle the Sensory-Coupling Interface}

The design of the modality interface determines how much relational structure is available for $\M$ to discover.
The literature offers three paradigms:

\paragraph{1. Linear/MLP projection (LLaVA, our current approach).}
Liu et al.~\cite{liu2024llava} showed that a simple 2-layer MLP mapping CLIP ViT outputs into the LLM's token space is sufficient for strong visual reasoning.
Image patches become ``visual tokens'' processed by the same attention mechanism as text tokens.
LLaVA-OneVision extends this to dynamic resolution (up to 7{,}290 visual tokens per image).
\textbf{For ORN:} this is the natural fit.
$\M$ operates on continuous geometry; linear projection preserves the geometric relationships $\M$ needs.
No information loss from discretisation.

\paragraph{2. Perceiver / Q-Former bottleneck (Flamingo, BLIP-2).}
Alayrac et al.\ (Flamingo) use a Perceiver Resampler: learned query tokens cross-attend to visual features, producing a fixed-length output (e.g.\ 64 tokens from 196 patches).
Li et al.\ (BLIP-2) use a Q-Former: 32 learned queries ``interview'' the visual patches.
Both act as information bottlenecks before the language model.
\textbf{For ORN:} likely redundant.
$\M$ with effective rank $r^*$ IS already a bottleneck---it selects the $r^*$ most important coupling dimensions.
Adding a separate bottleneck before $\M$ might discard relational structure that $\M$ could exploit.

\paragraph{3. Discrete tokenisation (VQ-VAE, VQGAN).}
Continuous signals are mapped to codebook indices, enabling uniform autoregressive generation across modalities.
A 2025 survey~\cite{discrete2025survey} confirms that discrete tokens still underperform continuous representations for \emph{understanding} tasks, though they enable generation.
\textbf{For ORN:} worst fit.
Discretisation destroys the continuous coupling geometry $\M$ operates on.

\paragraph{4. Native multimodal (Gemini, ImageBind).}
Gemini processes all modalities through shared transformer layers from the start---no separate encoders aligned post-hoc.
ImageBind~\cite{fang2023imagebind} creates a single shared embedding space for six modalities using only image-paired training data; transitive alignment through the image ``hub'' creates cross-modal alignment without explicit (audio, text) pairs.
\textbf{For ORN:} ImageBind's transitive alignment is structurally close to what $\M$ provides.
If $\M$ learns a shared coupling geometry, it should transitively align modalities that were never explicitly paired.

\subsection{The Neural ODE Connection}

The ORN's orbital iteration $h_{\ell+1} = h_\ell + f(\M, h_\ell)$ is a discretised Neural ODE~\cite{chen2018neural}: $dh/dt = f(\M, h(t))$.
$\M$ defines the vector field; the residual stream is the trajectory; each modality starts from a different initial condition.

This connects the ORN to a rich theoretical literature:
\begin{itemize}[nosep]
    \item Latent ODEs encode initial conditions and evolve them through continuous dynamics---exactly the orbital embedding idea (seeds + shared dynamics).
    \item The Neural Differential Manifold~\cite{ndm2025} treats network parameters as a metric tensor defining a Riemannian manifold---the closest existing work to $\M$-as-metric, though NDM computes a different metric per layer while ORN shares one.
    \item Hyperbolic embeddings~\cite{nickel2017poincare} represent hierarchical structure in curved space. If $\M$'s eigenspectrum creates a soft hierarchy (large eigenvalues for coarse coupling, small for fine), $\M$ implicitly induces hyperbolic-like geometry in coupling space.
\end{itemize}

\subsection{Design Principle for ORN Multimodal}

Based on the literature review AND the embedding ablation experiment below, the design principle is:
\begin{enumerate}[nosep]
    \item Use continuous projection (not discrete tokenisation)---$\M$ needs continuous geometry.
    \item Do not add a bottleneck before $\M$---$\M$ IS the bottleneck (rank $r^*$).
    \item Tell $\M$ what modality it is processing via a learnable modality token (e.g.\ \texttt{[IMG]}, \texttt{[AUD]}).  This is the single largest embedding improvement we found ($-10.0\%$ over transformer baseline at $d{=}96$).
    \item Normalise embeddings with RMSNorm before $\M$ sees them, so all modalities enter at consistent scale.  Combined with the modality token, this yields $-11.7\%$ vs transformer.
    \item Keep the projection simple.  MLP adds expressiveness but not performance; the gain is in \emph{what information reaches M}, not \emph{how non-linearly it is transformed}.
    \item Orbital modality embeddings (shared dynamics as alignment mechanism) do not help at current scale ($-2.1\%$ for $T{=}1$, $-3.1\%$ for $T{=}3$, both worse than the modality token alone).  The idea remains theoretically appealing but needs rethinking.
\end{enumerate}

\subsection{Embedding Ablation: How Should Modalities Enter M?}\label{sec:embedding-ablation}

To determine how modalities should be projected into $\M$'s coupling space, we test eight embedding variants on CIFAR-10 next-patch prediction (64 patches, $d{=}96$, $L{=}6$, 4 heads, 3000 steps, RMSNorm throughout).  All variants share the same SharedM backbone and response heads; only the input projection differs.  A transformer with standard per-layer $W_Q$, $W_K$ and linear embedding serves as baseline.

\begin{table}[h]
\centering
\small
\begin{tabular}{@{}lrrr@{}}
\toprule
\textbf{Variant} & \textbf{Params} & \textbf{MSE} & \textbf{vs TF} \\
\midrule
TF + Linear (baseline)      & 683,184  & 0.4486 & ---      \\
\midrule
SharedM + Linear             & 591,024  & 0.4366 & $-2.7\%$ \\
SharedM + MLP                & 604,976  & 0.4346 & $-3.1\%$ \\
SharedM + ModToken           & 591,120  & 0.4037 & $-10.0\%$ \\
SharedM + Normed             & 591,120  & 0.4319 & $-3.7\%$ \\
SharedM + Orbital($T{=}1$)  & 607,264  & 0.4391 & $-2.1\%$ \\
SharedM + Orbital($T{=}3$)  & 607,264  & 0.4346 & $-3.1\%$ \\
\textbf{SharedM + ModToken+Norm} & \textbf{591,216} & \textbf{0.3962} & $\mathbf{-11.7\%}$ \\
SharedM + ModToken+MLP      & 605,072  & 0.4115 & $-8.3\%$ \\
\bottomrule
\end{tabular}
\caption{Embedding ablation on CIFAR-10 next-patch prediction ($d{=}96$, 3000 steps, RMSNorm).  ModToken = learnable \texttt{[IMG]} token prepended to the patch sequence.  Norm = RMSNorm after projection.  All SharedM variants use identical backbone; only the embedding differs.  SharedM variants use ${\sim}13\%$ fewer parameters than the transformer baseline.}
\label{tab:embedding-ablation}
\end{table}

\paragraph{Key findings.}
\begin{enumerate}[nosep]
    \item \textbf{Modality tokens dominate.}  A single learnable \texttt{[IMG]} token prepended to the sequence gives $\M$ a routing signal: ``this is image data, apply image coupling rules.''  The $-10.0\%$ improvement over transformer baseline is the largest single-factor gain, achieved with only 96 extra parameters (the token itself).

    \item \textbf{RMSNorm + modality token stack.}  Normalising embeddings so all modalities enter at consistent scale ($-3.7\%$ alone) combines with the modality token for $-11.7\%$ total.  This is the recommended embedding for ORN multimodal.

    \item \textbf{MLP does not help.}  A 2-layer MLP (LLaVA 1.5 style) adds 14K parameters and delivers only $-3.1\%$ (same as orbital $T{=}3$).  When combined with a modality token, MLP is worse than simple linear+norm ($-8.3\%$ vs $-11.7\%$).  $\M$ does not need a rich embedding; it needs a well-scaled one with a modality signal.

    \item \textbf{Orbital embeddings underperform.}  Using $\M$'s own dynamics as an embedding function ($T{=}1$: $-2.1\%$; $T{=}3$: $-3.1\%$) helps less than a modality token alone ($-10.0\%$).  More iterations improve slightly but remain well below the modality token.  Orbital embeddings remain theoretically interesting (the dynamics as alignment mechanism) but are not practical at current scale.

    \item \textbf{All SharedM variants beat transformer.}  Even with ${\sim}13\%$ fewer parameters (591K vs 683K), every SharedM variant matches or exceeds transformer.
\end{enumerate}

\paragraph{Interpretation.}
$\M$ is a coupling authority, not an embedding function.  It works best when it receives inputs that are (a) well-scaled (RMSNorm) and (b) labelled by modality (modality token).  Given these two signals, $\M$ discovers the coupling geometry itself.  Enriching the projection (MLP) or folding M into the embedding process (orbital) adds complexity without payoff.  The message: prepare inputs simply, let $\M$ do the coupling.

% ======================================================================
\section{The Universal Coupling Manifold: 50K-Step Experiment}
\label{sec:universal-manifold}

The preceding experiments established that SharedM works across modalities at short training horizons (8K--20K steps).
This section reports a longer, more controlled experiment designed to test the \textbf{fiber bundle prediction}: that $\M$ discovers a modality-agnostic coupling manifold (the base space) through standard training alone, without any semigroup loss.

\subsection{Setup}

\textbf{Models.} ORN and Transformer, matched architecture ($d{=}256$, $L{=}8$, $H{=}8$), skew-dissipative $\M$ for the ORN.
The ORN has 6.0M parameters; the Transformer has 6.7M (11\% more, due to per-layer $W_Q, W_K$).

\textbf{Data.} All four modalities trained simultaneously via round-robin: CIFAR-10 image patches (50K$\times$64$\times$48), MAESTRO MIDI sequences (89K$\times$64$\times$6), ESC-50 environmental audio patches (1600$\times$96$\times$256), TinyStories text (9.5M bytes).

\textbf{Training.} Standard next-token/next-patch CE and MSE losses. \textbf{No semigroup loss.} \textbf{No CORN.} 50K steps, AdamW, $\text{lr} = 2 \times 10^{-4}$.

\textbf{Diagnostics} at checkpoints 10K, 25K, 50K:
\begin{itemize}[nosep]
    \item D1: coupling pattern PCA --- are attention patterns modality-agnostic?
    \item D2: frozen-$\M$ transfer --- freeze $\M$, train fresh response heads on text only
    \item D3: base-space orbit reconstruction --- can $\M^k h$ projected to the coupling manifold reconstruct content?
\end{itemize}

\subsection{Prediction Quality}

\begin{center}\small
\begin{tabular}{@{}lrrrrl@{}}
\toprule
& \textbf{Images} & \textbf{Music} & \textbf{Env audio} & \textbf{Text} & \textbf{Note} \\
\midrule
ORN (6.0M)        & 0.049 & 0.017 & \textbf{2.56} & 0.481 & 10\% fewer params \\
Transformer (6.7M) & \textbf{0.046} & \textbf{0.016} & 3.02 & \textbf{0.470} & \\
\bottomrule
\end{tabular}
\end{center}

The ORN is competitive on all modalities with 10\% fewer parameters.
It \textbf{wins on environmental audio} by 15\% (2.56 vs 3.02).
The transformer wins on images (+7\%) and text (+2\%).
Music is a tie.

Compared to the 20K-step results (Table above), the 50K run shows continued improvement on all modalities for both architectures: the ORN's environmental audio drops from 5.58 (20K) to 2.56 (50K), and text from 0.508 to 0.481.
Longer training helps, and the ORN does not saturate earlier than the transformer.

\subsection{D1: The Coupling Manifold Is Modality-Agnostic}

At each checkpoint, we extract attention patterns from images, music, and text, compute per-modality mean coupling features (entropy and max-attention statistics), and measure cross-modal cosine similarity.

\begin{center}\small
\begin{tabular}{@{}lrrr@{}}
\toprule
\textbf{Checkpoint} & \textbf{img $\leftrightarrow$ mus} & \textbf{img $\leftrightarrow$ txt} & \textbf{mus $\leftrightarrow$ txt} \\
\midrule
10K  & 0.9999 & 0.9987 & 0.9979 \\
25K  & 0.9999 & 0.9978 & 0.9967 \\
50K  & 0.9999 & 0.9978 & 0.9965 \\
\bottomrule
\end{tabular}
\end{center}

\textbf{Coupling patterns are $> 99.6\%$ similar across all modality pairs at all checkpoints.}
The coupling manifold is modality-agnostic from early training (10K) and remains so.
$\M$ does not develop modality-specific coupling modes---it uses the same geometric structure for images, music, environmental sounds, and text.

This is the multimodal analogue of the spectral invariance result for text transformers: the coupling rule is universal, and the per-layer variation comes from the evolving residual stream (the fiber), not from the coupling matrix (the base).

\subsection{D2: Frozen-M Transfer}

We freeze $\M$ from the multimodal ORN and train fresh response heads ($W_V, W_O$, FFN, corrections, embeddings) on text only for 3K steps.
The transfer gap measures how much worse the frozen-$\M$ model is compared to the original multimodal model on text.

\begin{center}\small
\begin{tabular}{@{}lrrr@{}}
\toprule
\textbf{Checkpoint} & \textbf{Original text loss} & \textbf{Transfer text loss} & \textbf{Gap} \\
\midrule
10K  & 0.512 & 0.508 & $-0.8\%$ (transfer \emph{better}) \\
25K  & 0.477 & 0.462 & $-3.0\%$ (transfer \emph{better}) \\
50K  & 0.481 & 0.491 & $+2.2\%$ \\
\bottomrule
\end{tabular}
\end{center}

At 50K, the frozen-$\M$ transfer costs only \textbf{2.2\%}.
$\M$ carries genuine geometric information that transfers to a new task---it is not an artefact of co-adaptation with the multimodal response heads.

At 10K and 25K, the transfer is \emph{better} than the original (negative gap).
This is because the fresh text-only response heads can fully specialise for text, unburdened by the multimodal interference that the original response heads must handle.
The base space (coupling geometry) is shared; the fiber (response heads) benefits from specialisation.

\subsection{Spectral Crystallisation}

$\M$'s effective rank over training:

\begin{center}\small
\begin{tabular}{@{}lr@{}}
\toprule
\textbf{Step} & \textbf{Effective rank} \\
\midrule
1 & 165 \\
10,000 & 36 \\
25,000 & 24 \\
50,000 & \textbf{17.5} \\
\bottomrule
\end{tabular}
\end{center}

The coupling manifold self-organises to $\sim$18 dimensions across four modalities.
This is consistent with the text-only experiments, where the base-space optimal dimension was $\sim$20 (confirmed by PCA of GPT-2 attention patterns at 8 dimensions for 95\% variance, and by the base-space orbit reconstruction experiments where $N_{\text{base}} = 20$ was optimal).

The rank stabilises between 25K and 50K (24 $\to$ 17.5), suggesting the manifold structure is largely determined by 25K steps.
The final rank of 17.5 means $\M$ uses fewer than 18 of its 256 available dimensions for coupling---a compression ratio of $14\times$.

\subsection{The Fiber Bundle Picture, Confirmed Multimodally}

These results confirm the central prediction of the fiber bundle framework:

\begin{enumerate}[nosep]
    \item \textbf{The base space is modality-agnostic.}
    $\M$'s coupling patterns are $> 99.6\%$ similar across images, music, audio, and text.
    The coupling manifold encodes \emph{how} things relate (structure), not \emph{what} they are (content).

    \item \textbf{The fiber encodes modality-specific content.}
    Frozen-$\M$ transfer costs only 2.2\%: the coupling geometry transfers; the response heads must re-specialise for each modality.
    At early checkpoints, fresh single-modality response heads outperform multimodal ones---the fiber benefits from specialisation, the base does not.

    \item \textbf{The manifold self-organises to $\sim$18 dimensions.}
    Despite having 256 dimensions available, $\M$ concentrates its coupling energy into 18 modes---across all four modalities simultaneously.
    This is the same dimensionality found in text-only experiments, suggesting the coupling manifold's intrinsic dimension is a property of relational structure itself, not of any specific modality.

    \item \textbf{No semigroup loss is needed.}
    Standard next-token/next-patch training discovers the shared manifold.
    The coupling geometry emerges from the structure of the task (``nearby things attend to each other''), not from an explicit composability constraint.
    CORN improves the manifold further (rank 17 $\to$ 11 at 20K steps), but the manifold exists without it.
\end{enumerate}

\begin{remark}
The orbit reconstruction diagnostic (D3) showed weak signal at 50K steps: text orbit accuracy 0--4\%, image MSE $\sim$0.28.
This used a simple mean-pool decoder without training---the C6b experiments (companion paper) show that a trained orbit decoder on the 20D base space achieves 86\% reconstruction accuracy.
The multimodal base-space orbits likely carry substantial information; the diagnostic was too weak to extract it.
A proper trained orbit decoder for multimodal base-space reconstruction is the natural next experiment.
\end{remark}

% ======================================================================
\section{The Orbit as Interlingua: Cross-Modal Convergence}
\label{sec:orbit-interlingua}

The shared coupling manifold (Section~\ref{sec:universal-manifold}) shows that $\M$'s coupling patterns are modality-agnostic ($> 99.6\%$ cosine similarity across all modality pairs).
But coupling similarity is not the same as \emph{embedding} similarity.
When we attempt standard multimodal tasks---zero-shot classification, cross-modal retrieval---the raw embeddings fail: CIFAR-10 zero-shot accuracy is 10.4\% (chance level), and cross-modal search precision is near zero.

The coupling geometry is shared, but the model cannot exploit it.
The input projections ($\text{img\_proj}$, $\text{txt\_emb}$) map different modalities into different regions of the $d$-dimensional space.
They share the same ``road network'' ($\M$'s coupling) but have no roads between them.

\subsection{Three Levels of Cross-Modal Similarity}

We probed the trained 50K-step multimodal ORN at three levels of representation, measuring cosine similarity between semantically matching content across modalities (e.g.\ dog images, dog bark audio, and the text ``dog'').
Four overlapping categories (airplane, cat, dog, frog) appear in both CIFAR-10 and ESC-50.

\begin{center}\small
\begin{tabular}{@{}lrrr@{}}
\toprule
\textbf{Representation level} & \textbf{img $\leftrightarrow$ txt} & \textbf{env $\leftrightarrow$ txt} & \textbf{img $\leftrightarrow$ env} \\
\midrule
Raw embedding (mean-pool $h$) & 0.62 & 0.46 & 0.71 \\
Base-space projection ($V_k^\top h$) & 0.93 & 0.78 & 0.90 \\
Orbit $k{=}4$ ($\M_{\text{base}}^4 V_k^\top h$) & 0.99 & 0.98 & 1.00 \\
Orbit $k{=}8$ & \textbf{1.000} & \textbf{1.000} & \textbf{1.000} \\
\bottomrule
\end{tabular}
\end{center}

\textbf{$\M$'s orbit dynamics are a convergent funnel.}
All modalities---images, environmental audio, and text---converge to cosine similarity \textbf{1.000} in the base space after $\sim$4--5 orbit steps.
The orbit IS the interlingua.

The progression reveals the geometry:
\begin{itemize}[nosep]
    \item \textbf{Raw embeddings} (0.46--0.71): the input projections place different modalities in different regions.
    The shared coupling cannot help if the addresses are disjoint.
    \item \textbf{Base-space projection} (0.78--0.93): projecting onto $\M$'s top 20 eigenvectors already recovers substantial cross-modal alignment.
    The coupling manifold is more aligned than the full hidden state because the fiber (modality-specific content) is discarded.
    \item \textbf{Orbit iteration} ($\to$ 1.000): $\M$'s dynamics systematically collapse the remaining modality gap.
    Each orbit step amplifies the dominant eigenmode and suppresses the rest.
    By $k{=}4$, cross-modal similarity exceeds 0.98.
    By $k{=}8$, it reaches 1.000.
\end{itemize}

\subsection{The Convergence Is Too Aggressive}

There is a critical subtlety: the orbit converges to cosine 1.000 for \emph{all} inputs, not just semantically matching ones.
The orbit of a dog image and the orbit of an airplane image also converge to cosine $\sim$1.000---because $\M$'s dominant eigenmode absorbs everything.

The L2 separation tells a different story.
At the orbit endpoint ($k{=}8$), the L2 distance between same-class cross-modal pairs is $2.1\times$ smaller than between different-class pairs (for img$\leftrightarrow$txt) and $2.6\times$ smaller (for img$\leftrightarrow$env).
The class information is encoded in the \emph{magnitude profile} across eigenmodes, not in the final direction.

\begin{center}\small
\begin{tabular}{@{}lrrl@{}}
\toprule
\textbf{Distance type} & \textbf{L2 dist} & \textbf{Ratio} & \textbf{Meaning} \\
\midrule
Same-class img$\leftrightarrow$txt & $d_{\text{same}}$ & 1.0$\times$ & Baseline \\
Same-class img$\leftrightarrow$env & $d_{\text{same}}$ & 1.0$\times$ & \\
Diff-class img$\leftrightarrow$img & $d_{\text{diff}}$ & \textbf{2.1--2.6$\times$} & Same-class is closer \\
\bottomrule
\end{tabular}
\end{center}

\textbf{The orbit provides modality-agnostic direction but not class-discriminative direction.}
Everything converges to $\M$'s leading eigenmode.
The discriminative signal is in how fast each eigenmode decays---the transient dynamics, not the fixed point.

\subsection{The Interlingua Hypothesis}

This suggests a precise training strategy.
The orbit is not the final representation---it is the \emph{regulariser} that shows the encoder where the shared geometry wants things to land.
The useful cross-modal representation is at \textbf{early orbit steps} ($k{=}1$--3), where there is both:
\begin{itemize}[nosep]
    \item Cross-modal alignment (cosine 0.85--0.96, much better than raw embeddings)
    \item Class discrimination (the eigenmodes have not yet collapsed to the dominant one)
\end{itemize}

The training principle: \textbf{train through the orbit}.
Use the orbit at $k{=}2$--3 as the space where cross-modal matching happens.
The encoder learns to place semantically similar content from different modalities at the same base-space location, guided by $\M$'s dynamics.

\subsection{Proposed: Orbit-Mediated Cross-Modal Training}

Three training objectives, to be tested individually and in combination:

\paragraph{O1: Orbit contrastive alignment.}
For semantically matched pairs (e.g.\ image patches of a cat and the text ``cat''), compute orbit points at $k{=}2$ in the base space.
Apply an InfoNCE contrastive loss:
\begin{equation}
    \mathcal{L}_{\text{O1}} = -\log \frac{\exp(\text{sim}(\text{orb}_k^{\text{img}}, \text{orb}_k^{\text{txt}}) / \tau)}{\sum_{j} \exp(\text{sim}(\text{orb}_k^{\text{img}}, \text{orb}_k^{\text{txt}_j}) / \tau)}
\end{equation}
where $\text{orb}_k = \M_{\text{base}}^k V_k^\top h$ and $\text{sim}$ is cosine similarity.

This is CLIP-like, but the contrastive matching happens in the orbit space, not the raw embedding space.
The orbit provides the inductive bias: the encoder only needs to get representations ``close enough'' in base space for $\M$'s dynamics to pull them together.
This should be easier to learn than raw embedding alignment.

\paragraph{O2: Cross-modal orbit prediction.}
Given an image, predict the orbit of the corresponding text (and vice versa):
\begin{equation}
    \mathcal{L}_{\text{O2}} = \| \text{orb}_k^{\text{img}} - \text{sg}(\text{orb}_k^{\text{txt}}) \|^2 + \| \text{sg}(\text{orb}_k^{\text{img}}) - \text{orb}_k^{\text{txt}} \|^2
\end{equation}
where $\text{sg}(\cdot)$ is stop-gradient (BYOL-style, to prevent collapse).
This trains the encoders to produce orbits that match across modalities, without negative pairs.

\paragraph{O3: Orbit-grounded next-token prediction.}
During standard next-token/next-patch training, randomly replace the input modality while keeping the prediction target.
Feed image patches of a cat, but the loss is on predicting the \emph{text} ``a small furry cat sat on the mat.''
The representation must bridge modalities to predict correctly.

Crucially, the cross-modal prediction goes through the orbit: the image is encoded, projected to base space, evolved $k$ steps by $\M_{\text{base}}$, lifted back to $d$ dimensions, and decoded by the text response head.
The orbit is the bottleneck that forces the representation through the shared geometry.

\paragraph{Training protocol.}
\begin{enumerate}[nosep]
    \item Phase 1 (10K steps): standard multimodal next-token/next-patch training (build good per-modality representations).
    \item Phase 2 (10K steps): add O1 (orbit contrastive) at $\lambda{=}0.1$.
    The orbit provides the cross-modal bridge.
    \item Phase 3 (10K steps): add O3 (orbit-grounded cross-modal prediction) at $\lambda{=}0.1$.
    The model learns to predict across modalities through the orbit.
    \item Evaluate: zero-shot classification, cross-modal retrieval, orbit alignment probe.
\end{enumerate}

The key comparison: does orbit-mediated contrastive learning outperform standard (raw embedding) contrastive learning?
If yes, $\M$'s geometry provides a genuine inductive bias for cross-modal alignment---not just shared coupling patterns, but a usable bridge between modalities.

\paragraph{What this would prove.}
If O1 or O3 produces zero-shot classification accuracy significantly above chance (${>}30\%$ on CIFAR-10, ${>}5\%$ on ESC-50) with modest training ($\sim$30K steps at $d{=}256$), then the ORN offers a genuine multimodal advantage: the shared coupling manifold provides geometric structure that makes cross-modal alignment \emph{easier to learn} than in a standard transformer.
The orbit is the mechanism: it funnels different modalities toward the same region of base space, and the encoder only needs to learn a rough alignment for $\M$'s dynamics to refine it.

If it fails, the shared coupling geometry is a coincidence of training dynamics, not a usable bridge---and the ORN's multimodal story reduces to parameter efficiency alone.

\section{What Would Confirm or Falsify}

\paragraph{Confirms the theory:}
\begin{itemize}[nosep]
    \item Spectral invariance holds across modalities (Exp M1)
    \item SharedM matches/beats transformer on image-text (Exp M2) \checkmark
    \item $r^*$ increases modestly with added modalities (Exp M3)
    \item Same architecture works for audio without modification (Exp M4)
    \item Sub-linear rank growth confirmed at $d{=}96$: 3 modalities need rank 3 vs sum 10 (Table~\ref{tab:confirmed-rank}) \checkmark
    \item Cross-modal transfer works: text $\M$ transfers to images with $+3.2\%$ penalty (Table~\ref{tab:frozen-transfer}) \checkmark
    \item Music coupling is rank 2, confirming very low-dimensional coupling geometry (Table~\ref{tab:maestro}) \checkmark
    \item SharedM beats transformer on music ($-3.7\%$) and images ($-0.7\%$) independently \checkmark
    \item Crossover trajectory confirmed: ORN trails transformer at 3K steps but overtakes on images ($-2.6\%$) and environmental sounds ($-4.2\%$) by 8K steps \checkmark
    \item Coupling manifold is modality-agnostic: $> 99.6\%$ cosine similarity across all 4 modality pairs at 50K steps (Section~\ref{sec:universal-manifold}) \checkmark
    \item Frozen-$\M$ transfer costs only $2.2\%$ from 4-modality training to single-modality text (Section~\ref{sec:universal-manifold}) \checkmark
    \item Coupling manifold self-organises to $\sim$18 dimensions across 4 modalities, matching text-only result (Section~\ref{sec:universal-manifold}) \checkmark
    \item No semigroup loss needed: standard training discovers the shared manifold (Section~\ref{sec:universal-manifold}) \checkmark
\end{itemize}

\paragraph{Falsifies the theory:}
\begin{itemize}[nosep]
    \item Cross-modal coupling has low spectral correlation ($< 0.7$) --- \textbf{REJECTED}: measured $> 0.996$ \checkmark
    \item SharedM degrades significantly ($> 10\%$) on multimodal tasks --- \textbf{REJECTED}: ORN wins on env audio ($-15\%$), competitive elsewhere \checkmark
    \item $r^*$ doubles or more with each added modality --- \textbf{REJECTED}: rank 17.5 for 4 modalities vs $\sim$20 for text alone \checkmark
    \item If frozen-M test shows $> 10\%$ worse than from-scratch --- \textbf{REJECTED}: gap is $+2.2\%$ \checkmark
\end{itemize}

All four falsification criteria have been tested and rejected.
The theory survives every test we have designed for it.

% ======================================================================

% ======================================================================
\section{The Emerging Consensus: Coupling Is Shared, Response Is Not}
\label{sec:consensus}
% ======================================================================

A striking convergence has emerged in the 2024--2025 multimodal literature: across diverse architectures and scales, the \emph{same decomposition} appears independently.
Attention coupling (which tokens relate to which) can be shared across modalities; response mechanisms (FFN, projections) benefit from being modality-specific.
This is precisely the ORN's coupling--response decomposition---arrived at from a different direction.

\paragraph{The Platonic Representation Hypothesis.}
Huh et al.~\cite{huh2024platonic} demonstrate that as vision and language models scale, their internal representations \emph{converge}: different architectures trained on different objectives and different modalities increasingly measure distance between datapoints in the same way.
The convergence is not engineered---it emerges from scale and task pressure.
They argue this reflects a shared ``platonic representation'' of reality's statistical structure.
\textbf{For ORN:} if representations converge across modalities and architectures, then the coupling geometry (what relates to what) is genuinely modality-invariant.
A shared $\M$ that captures this universal relational structure is exactly what the Platonic Representation Hypothesis predicts should exist.
Our measurement of $>99.6\%$ coupling similarity across 4 modalities (Section~\ref{sec:universal}) is direct evidence for this convergence at the coupling level.

\paragraph{Meta-Transformer: frozen encoder across 12 modalities.}
Zhang et al.~\cite{zhang2023metatransformer} demonstrate that a single \emph{frozen} ViT encoder---pretrained on one modality---processes 12 modalities (text, image, point cloud, audio, video, time series, tabular, X-ray, infrared, hyperspectral, graph, IMU) without any fine-tuning of the encoder.
Only lightweight modality-specific tokenisers and task heads are updated.
\textbf{For ORN:} Meta-Transformer's frozen encoder is the closest existing architecture to shared $\M$.
The entire attention mechanism (Q, K, V, FFN) learned from one modality transfers to 11 others.
This validates the strongest version of our claim: coupling geometry is modality-universal.

\paragraph{Mixture-of-Transformers: shared attention, separate response.}
Liang et al.~\cite{liang2024mot} decouple FFN, attention projections, and LayerNorm by modality while keeping global self-attention shared across all modalities.
This matches the Chameleon 7B baseline at 55.8\% of FLOPs for text-image, and 37.2\% for text-image-speech.
\textbf{For ORN:} MoT's architecture \emph{independently discovers} the ORN decomposition.
Share the coupling rule (global attention), keep the response modality-specific (FFN, projections).
The key difference: MoT uses per-layer shared attention; ORN memoise across layers too.

\paragraph{ONE-PEACE: shared self-attention for vision, audio, and language.}
Wang et al.~\cite{wang2023onepeace} combine shared self-attention layers with modality adapters and modality-specific FFNs.
Two modality-agnostic pretraining objectives (cross-modal aligning contrast + intra-modal denoising contrast) align semantic spaces across vision, audio, and language at 4B parameters.
\textbf{For ORN:} ONE-PEACE demonstrates at production scale that shared self-attention + modality-specific FFN works.
The intra-modal denoising objective connects to our phased training harness (Phase 2: fiber denoising under frozen $\M$).

\paragraph{4M-21: one transformer for 21 modalities.}
Bachmann et al.~\cite{bachmann2024fourm} train a single 3B-parameter transformer encoder-decoder on 21 diverse modalities (RGB, depth, normals, semantics, edges, DINOv2, SAM labels, metadata, etc.) via masked modelling over discrete tokens.
\textbf{For ORN:} 4M-21 is the strongest existence proof that shared attention works across many modalities.
At 21 modalities and 3B parameters, the shared coupling structure has not collapsed.
The tokenisers handle modality-specific encoding; the transformer's attention is the universal coupling rule.

\paragraph{Dual Attention Transformer: relational vs sensory information.}
Altabaa \& Lafferty~\cite{altabaa2024dat} factor standard attention into two types: \emph{sensory heads} (properties of individual tokens) and \emph{relational heads} (relationships between tokens).
The relational component generalises across language, vision, and symbolic reasoning.
\textbf{For ORN:} DAT's relational heads are conceptually identical to shared $\M$---they isolate the coupling rule from the content response.
The key difference: DAT uses per-layer relational heads; ORN memoises the relational structure into a single shared $\M$.

\paragraph{Synthesis: the ORN as explicit architecture for an emergent consensus.}
These six independent lines of work converge on the same factorisation:
\begin{center}\small
\begin{tabular}{@{}lll@{}}
\toprule
\textbf{Paper} & \textbf{Shared across modalities} & \textbf{Modality-specific} \\
\midrule
Meta-Transformer~\cite{zhang2023metatransformer} & Frozen encoder (all attention + FFN) & Tokenisers, task heads \\
MoT~\cite{liang2024mot} & Global self-attention & FFN, projections, LayerNorm \\
ONE-PEACE~\cite{wang2023onepeace} & Self-attention layers & Adapters, FFN \\
4M-21~\cite{bachmann2024fourm} & Entire transformer & Tokenisers \\
DAT~\cite{altabaa2024dat} & Relational attention heads & Sensory heads \\
\midrule
\textbf{ORN V3} & \textbf{$\M$ + wide shared FFN} & \textbf{Response heads, corrections} \\
\bottomrule
\end{tabular}
\end{center}

The ORN makes this consensus \emph{explicit and measurable}.
Rather than treating shared attention as an engineering convenience, we track the spectral geometry of the shared coupling tensor $\M$: its effective rank (how many coupling dimensions are needed), its eigenmode partitioning (which modes serve which modalities), and its crystallisation dynamics (when does the shared geometry stabilise).
The sub-linear rank growth measured in Section~\ref{sec:sublinear-rank} (3 modalities need rank 3, not rank 10) is a quantitative prediction that none of the above papers make but all would benefit from testing.

\paragraph{V3: shared FFN as a second axis of universality.}
ORN V3 extends the sharing from coupling ($\M$) to content (wide shared FFN).
This is closer to Meta-Transformer (which shares everything) and farther from MoT/ONE-PEACE (which keep FFN separate).
The experimental validation (Section~\ref{sec:v3-multimodal-geometry}) tests whether the shared FFN creates a cross-modal bridge that V2's per-layer FFN could not: a common representation space where all modalities meet in the same key-value memory.

\paragraph{What ORN uniquely contributes.}
No existing work measures the \emph{geometry} of the shared coupling.
The Platonic Representation Hypothesis predicts convergence; the ORN provides a way to measure it.
$\M$'s eigenspectrum is the coupling manifold made tangible: a finite-dimensional space whose rank, eigenmode partitioning, and crystallisation dynamics reveal what relational structure is shared across modalities and what is modality-specific.
The proposed V3 multimodal geometry experiment (Section~\ref{sec:v3-multimodal-geometry}) directly measures eigenmode usage per modality---which of $\M$'s singular vectors each modality activates---providing the first empirical map of the universal coupling manifold.


% ======================================================================
\section{V3 Multimodal Geometry Experiment}
\label{sec:v3-multimodal-geometry}
% ======================================================================

The key question is not whether multimodal works (it does---our coupling similarity is $>99.6\%$), but \emph{how much of the coupling geometry is intrinsically shared vs modality-specific}, and whether V3's shared FFN creates the cross-modal bridge that V2 lacked.

\subsection{How Literature Informed Each Design Choice}

Every design decision in this experiment is motivated by a specific finding in the literature or our own prior failures.
The goal is not to ``hack on multimodal'' but to measure what is fundamentally shared in the coupling space.

\paragraph{1. The tokeniser is everything (Meta-Transformer~\cite{zhang2023metatransformer}).}
Meta-Transformer shares a \emph{frozen} encoder across 12 modalities---and it works.
The only modality-specific component is the tokeniser.
This told us to stop trying to fix $\M$ or the backbone and instead fix the \emph{input interface}.
Our prior failures ($0\%$ cross-modal retrieval, Section~\ref{sec:universal}) happened with separate input projections: each modality mapped to a different region of $\R^d$.
$\M$'s coupling geometry was shared ($>99.6\%$ similarity) but the modalities never met in representation space.
Variant~B addresses this directly: a shared projection forces all modalities into the same region, so $\M$ can couple them.

\paragraph{2. Coupling is shareable, FFN is not---or is it? (MoT~\cite{liang2024mot}).}
MoT shares attention but separates FFN per modality, achieving 55\% FLOP savings.
ORN V3 does the opposite: it shares \emph{both} $\M$ and FFN.
This creates a direct test.
MoT predicts that shared FFN should hurt cross-modal performance.
Our wide shared FFN hypothesis predicts that one common key-value memory \emph{forces} cross-modal patterns into the same representation, creating a bridge that separate FFNs cannot.
We measure this via CKA alignment: if V3's shared FFN shows higher cross-modal CKA than Variant~A (separate projections, analogous to MoT's decomposition), shared FFN is creating the bridge.

\paragraph{3. Measure the kernel, not the embeddings (Platonic Representation Hypothesis~\cite{huh2024platonic}).}
Huh et al.\ show convergence at the level of the \emph{distance structure} (kernel), not at the level of raw representations.
$\M = AB^\top$ defines exactly such a coupling metric.
This told us to measure \textbf{eigenmode usage per modality}---which modes of $\M$'s singular value decomposition each modality activates---rather than raw embedding cosine similarity (which is low by construction for different modalities).
If the top-$k$ eigenmodes are shared, the coupling \emph{metric} has converged, exactly as Platonic predicts.

\paragraph{4. Asymmetric relations outperform symmetric (DAT~\cite{altabaa2024dat}).}
Altabaa \& Lafferty factor attention into relational (structural) and sensory (content) heads, finding that \emph{asymmetric} relational attention outperforms symmetric---directly paralleling our finding that $AB^\top$ beats $LL^\top$.
This validates tracking $\M$'s asymmetry ($\|\M - \M^\top\|/\|\M\|$) throughout training.
If asymmetry is consistent across modalities, the relational geometry is genuinely universal.

\paragraph{5. Middle layers are naturally modality-invariant (Uni-X~\cite{uni2025x}).}
Gradient conflict analysis shows that shallow and deep layers carry modality-specific statistics, but middle layers develop modality-invariant representations.
This told us to measure CKA \textbf{at each layer separately}.
We predict a U-shape: low CKA at layer~1 (modality-specific input), peak at layers~3--4 (middle, where $\M$'s coupling dominates), dropping at layer~6 (modality-specific output).
Observing this U-shape would confirm that $\M$ performs the modality-invariant coupling in the middle of the network, with per-layer response heads handling modality-specific work at the edges.

\paragraph{6. Crystallise before specialising (our frozen-M result $+$ developmental biology).}
Our frozen-M transfer gap ($+0.73\%$) plus Meta-Transformer's frozen encoder show that coupling geometry should be established before modality-specific training.
Variant~C implements this: a dedicated $\M$ crystallisation phase where $\M$ sees all modalities in alternating steps before response heads specialise.
This mirrors biological development (connectivity before specialisation---realisation~114 from the research diary) and the phased training harness designed on Day~6.

\subsection{Design}

Three training strategies on V3 ($d{=}128$, $L{=}6$, 6x wide shared FFN, three modalities: CIFAR-10 patches + TinyStories text + synthetic audio spectrograms):

\begin{description}[nosep]
\item[A: Separate projections (baseline).] Each modality has its own linear projection into $d$-space.
This is what failed before: fibers are disjoint, coupling similarity $>99.6\%$ but zero cross-modal retrieval.
Included as control---Meta-Transformer's lesson predicts this should underperform.
\item[B: Shared projection + modality tokens.] All modalities project through the \emph{same} linear layer into $d$-space.
A learned modality token (e.g.\ \texttt{[IMG]}, \texttt{[AUD]}, \texttt{[TXT]}) prepended to each sequence tells the model which modality it processes.
Forces all modalities into the same region of representation space.
The shared projection ensures all modalities enter $\M$'s coupling space through the same door; the modality token provides just enough information for the response heads to adapt.
This is the minimal intervention that Meta-Transformer's success predicts should work.
\item[C: B + M crystallisation phase.] Same as B, but before joint training, run 2000 steps of alternating all-params / M-only training.
$\M$-only steps (every 3rd step) freeze everything except $A$ and $B$, forcing $\M$ to discover coupling geometry that works for \emph{all} modalities simultaneously.
The response heads cannot co-adapt during these steps, so $\M$ must encode a genuinely universal structure.
\end{description}

\subsection{Measurements}

We measure geometry, not benchmarks:

\begin{enumerate}[nosep]
\item \textbf{Eigenmode usage per modality.} For each modality, project middle-layer representations onto $\M$'s singular vectors, measure energy per mode. If modalities share the top-$k$ modes, the coupling geometry is genuinely universal.
\item \textbf{CKA alignment at each layer.} Linear centered kernel alignment between modality pairs at every layer. Tests whether representations converge with depth (Uni-X~\cite{uni2025x} predicts middle layers are most aligned).
\item \textbf{Coupling similarity.} Direct comparison of $Q K^\top$ attention patterns across modalities.
\item \textbf{Eigenmode overlap.} Of the top-20 most-used modes per modality, how many are shared by all three?
\item \textbf{Frozen-backbone transfer.} Freeze $\M$ + shared FFN (the entire V3 backbone), train only per-layer response heads on a single modality. Measures how much multimodal geometry transfers.
\end{enumerate}

\subsection{Predictions}

If the Platonic Representation Hypothesis holds at the coupling level:
\begin{itemize}[nosep]
\item Variant B should show higher eigenmode overlap than A (shared projection forces co-habitation).
\item Variant C should show higher CKA alignment in middle layers (crystallisation before specialisation).
\item The top-10 eigenmodes should be $>50\%$ shared across all three modalities.
\item Frozen-backbone transfer should degrade $<5\%$ per modality.
\end{itemize}

If the shared FFN creates a cross-modal bridge (the V3 hypothesis):
\begin{itemize}[nosep]
\item CKA alignment should be higher than in equivalent V2 experiments (per-layer FFN).
\item The shared FFN's activation patterns should show modality-mixed usage (not partitioned by modality).
\end{itemize}


\section{Theory-Aligned Multimodal: From 0\% to 60\%}
\label{sec:theory-aligned-results}

\subsection{The Failure and the Fix}

Our first multimodal experiment (V1) achieved 0\% cross-modal retrieval despite $>99.6\%$ coupling similarity across modalities.
Diagnosis: three design bugs broke $\M$-alignment:
\begin{enumerate}[nosep]
    \item \textbf{Separate modality tokens} --- different learned starting points for text and image guaranteed different orbits regardless of $\M$.
    \item \textbf{Unpaired contrastive} --- random text pushed toward random images. InfoNCE without semantic pairing is noise.
    \item \textbf{No input distribution alignment} --- the text embedding and image projection produced different output distributions. $\M$ could trivially distinguish modalities from input statistics alone.
\end{enumerate}

The theory requires: $\M$ must be modality-agnostic.
Both modalities must enter $\M$'s coupling space in the \emph{same} distribution.
Three fixes, each derived from this principle:
\begin{enumerate}[nosep]
    \item \textbf{Single shared CLS token} --- both modalities start from the same point. Modality information comes from content, not from the starting token.
    \item \textbf{Shared input normalisation} (RMSNorm after projection) --- forces the same distribution into $\M$.
    \item \textbf{Paired contrastive} from CIFAR-10 labels --- pseudo-captions (``a picture of a dog'') paired with their images via InfoNCE.
\end{enumerate}

\subsection{Results: V2 Theory-Aligned Design}

Trained from scratch: $d{=}256$, $L{=}8$, $6\times$ wide shared FFN, 31M parameters, 28K steps on MPS ($\sim$6 hours).
Joint text (TinyStories) + image (CIFAR-10) training with paired contrastive from step 1.

\begin{center}\small
\begin{tabular}{@{}rrrrr@{}}
\toprule
Step & CIFAR-10 & Cross-modal & $\M$ rank & SV[0] \\
\midrule
2,000 & 29.6\% & 29.4\% & 180 & 2.59 \\
4,000 & 33.4\% & 33.1\% & 170 & 3.43 \\
8,000 & 46.1\% & 41.7\% & 160 & 4.46 \\
12,000 & 54.1\% & 50.2\% & 149 & 5.70 \\
16,000 & 60.1\% & 55.6\% & 142 & 6.42 \\
20,000 & 63.1\% & 59.7\% & 141 & 6.55 \\
\textbf{28,000} & \textbf{64.2\%} & \textbf{60.5\%} & \textbf{141} & \textbf{6.55} \\
\midrule
Chance & 10\% & 10\% & --- & --- \\
\bottomrule
\end{tabular}
\end{center}

Cross-modal retrieval: \textbf{60.5\%} (6$\times$ chance).
$\M$ crystallises from rank 180 $\to$ 141 while training on \emph{both} modalities simultaneously.
SV[0] grows monotonically (2.59 $\to$ 6.55), indicating $\M$ is sharpening its coupling geometry.

\subsection{What This Confirms}

\begin{enumerate}[nosep]
    \item \textbf{$\M$ bridges modalities when the input interface is designed correctly.}
    The same $\M$ that was stuck at 0\% in V1 reaches 60.5\% in V2.
    The coupling geometry was never the problem --- the input projections were.

    \item \textbf{$\M$ crystallises on multiple modalities simultaneously.}
    Rank 180 $\to$ 141 is normal crystallisation, comparable to text-only training (V2: 332 $\to$ 69 at $d{=}576$).
    The coupling geometry accommodates both text and images without conflict.

    \item \textbf{Pseudo-captions from class labels suffice as paired contrastive signal.}
    No image-caption dataset needed.
    Simple ``a picture of a [class]'' labels provide enough semantic pairing for InfoNCE to bridge modalities.
    Real image-text pairs (COCO, Conceptual Captions) would presumably yield stronger alignment.

    \item \textbf{The contrastive loss continues to drop throughout training} (2.08 $\to$ 0.30), indicating the modalities are still converging at step 28K.
    Longer training or stronger contrastive signal could push retrieval higher.
\end{enumerate}

\subsection{V1 vs V2 Comparison}

\begin{center}\small
\begin{tabular}{@{}lrr@{}}
\toprule
& V1 (broken) & V2 (theory-aligned) \\
\midrule
Cross-modal retrieval & $\approx 0\%$ & \textbf{60.5\%} \\
CIFAR-10 & 47\% & \textbf{64.2\%} \\
CLS token & Per-modality & Shared \\
Input normalisation & None & Shared RMSNorm \\
Contrastive data & Unpaired & Paired (class labels) \\
\bottomrule
\end{tabular}
\end{center}

Same architecture, same $\M$, same training budget.
The difference is three design principles, each following from the requirement that $\M$ must be modality-agnostic.

\begin{thebibliography}{99}

\bibitem{dosovitskiy2021vit}
A.~Dosovitskiy et al.
\newblock An Image is Worth 16x16 Words: Transformers for Image Recognition at Scale.
\newblock \emph{ICLR}, 2021.

\bibitem{radford2021clip}
A.~Radford et al.
\newblock Learning Transferable Visual Models From Natural Language Supervision.
\newblock \emph{ICML}, 2021.

\bibitem{gu2023mamba}
A.~Gu and T.~Dao.
\newblock Mamba: Linear-Time Sequence Modeling with Selective State Spaces.
\newblock \emph{arXiv:2312.00752}, 2023.

\bibitem{chen2020igpt}
M.~Chen, A.~Radford, R.~Child, J.~Wu, H.~Jun, D.~Luan, and I.~Sutskever.
\newblock Generative Pretraining from Pixels.
\newblock \emph{ICML}, 2020.

\bibitem{he2022mae}
K.~He, X.~Chen, S.~Xie, Y.~Li, P.~Doll\'{a}r, and R.~Girshick.
\newblock Masked Autoencoders Are Scalable Vision Learners.
\newblock \emph{CVPR}, 2022.

\bibitem{tian2024var}
K.~Tian, Y.~Jiang, Z.~Yuan, B.~Peng, and L.~Wang.
\newblock Visual Autoregressive Modeling: Scalable Image Generation via Next-Scale Prediction.
\newblock \emph{NeurIPS}, 2024. (Best Paper Award)

\bibitem{li2024continuous}
T.~Li et al.
\newblock Autoregressive Image Generation without Vector Quantization.
\newblock \emph{arXiv:2406.11838}, 2024.

\bibitem{friston2010free}
K.~Friston.
\newblock The free-energy principle: a unified brain theory?
\newblock \emph{Nature Reviews Neuroscience}, 11(2):127--138, 2010.

\bibitem{damasio2009cdz}
A.~Damasio.
\newblock Convergence and divergence in a neural architecture for recognition and memory.
\newblock \emph{Trends in Neurosciences}, 2009.

\bibitem{fang2023imagebind}
Y.~Girdhar et al.
\newblock ImageBind: One Embedding Space To Bind Them All.
\newblock \emph{CVPR}, 2023.

\bibitem{liu2024llava}
H.~Liu et al.
\newblock Visual Instruction Tuning.
\newblock \emph{NeurIPS}, 2024.

\bibitem{discrete2025survey}
Discrete Tokenization for Multimodal LLMs: A Comprehensive Survey.
\newblock \emph{arXiv:2507.22920}, 2025.

\bibitem{nickel2017poincare}
M.~Nickel and D.~Kiela.
\newblock Poincar\'{e} Embeddings for Learning Hierarchical Representations.
\newblock \emph{NeurIPS}, 2017.

\bibitem{ndm2025}
The Neural Differential Manifold.
\newblock \emph{arXiv:2510.25113}, 2025.

\bibitem{chen2018neural}
R.~T.~Q.~Chen, Y.~Rubanova, J.~Bettencourt, and D.~Duvenaud.
\newblock Neural Ordinary Differential Equations.
\newblock \emph{NeurIPS}, 2018.

\bibitem{huh2024platonic}
M.~Huh, B.~Cheung, T.~Wang, and P.~Isola.
\newblock The Platonic Representation Hypothesis.
\newblock \emph{ICML}, 2024.

\bibitem{zhang2023metatransformer}
Y.~Zhang, K.~Gong, K.~Zhang, H.~Li, Y.~Qiao, W.~Ouyang, and X.~Yue.
\newblock Meta-Transformer: A Unified Framework for Multimodal Learning.
\newblock \emph{ICLR}, 2024.

\bibitem{liang2024mot}
W.~Liang, L.~Yu, L.~Luo, S.~Iyer, N.~Dong, C.~Zhou, G.~Ghosh, M.~Lewis, W.-t.~Yih, L.~Zettlemoyer, and X.~V.~Lin.
\newblock Mixture-of-Transformers: A Sparse and Scalable Architecture for Multi-Modal Foundation Models.
\newblock \emph{arXiv:2411.04996}, 2024.

\bibitem{wang2023onepeace}
P.~Wang et al.
\newblock ONE-PEACE: Exploring One General Representation Model Toward Unlimited Modalities.
\newblock \emph{NeurIPS}, 2023.

\bibitem{bachmann2024fourm}
R.~Bachmann, O.~F.~Kar, D.~Mizrahi, A.~Garjani, M.~Gao, D.~Griffiths, J.~Hu, A.~Dehghan, and A.~Zamir.
\newblock 4M-21: An Any-to-Any Vision Model for Tens of Tasks and Modalities.
\newblock \emph{NeurIPS}, 2024.

\bibitem{altabaa2024dat}
A.~Altabaa and J.~Lafferty.
\newblock Disentangling and Integrating Relational and Sensory Information in Transformer Architectures.
\newblock \emph{ICML}, 2025.

\bibitem{uni2025x}
Uni-X: Resolving Gradient Conflicts in Unified Multimodal Transformers.
\newblock \emph{arXiv:2509.24365}, 2025.

\end{thebibliography}


% ======================================================================
\section{V3 GPU Failure and V4 Corrected Design}
\label{sec:v3-failure}

The V3 multimodal run (605M, RTX 4090, Vast.ai, 10 days) failed.
Two bugs compounded to produce zero image learning and degrading text quality.
Both were found by targeted MPS experiments after the GPU run was killed.

\subsection{Bug 1: CLS Causal Blindness (Primary)}

The V2 theory-aligned design (Section~\ref{sec:theory-aligned-results}) used a shared CLS token with bidirectional attention for all modalities.
The V3 GPU implementation changed the attention pattern: \texttt{is\_causal=True} for text encoding, \texttt{is\_causal=False} for images.
This was intended to maintain autoregressive structure for text while allowing bidirectional for images.

The consequence was catastrophic: with causal masking, the CLS token at position 0 can only attend to \emph{itself}.
It cannot see any text content at positions $1, 2, \ldots, S$.
The contrastive loss was aligning image CLS (bidirectional, full visibility) with text CLS (causal, blind to text).
There was no gradient signal. Contrastive loss was stuck at $\ln(\text{batch\_size})$ for the entire run.

\paragraph{Confirmation.}
The pooling diagnostic experiment tested four pooling strategies on a matched 7M model trained for 3K steps:

\begin{center}\small
\begin{tabular}{lrrr}
\toprule
\textbf{Pooling} & \textbf{R@1} & \textbf{R@5} & \textbf{Text Val} \\
\midrule
CLS (broken) & 0.005 & 0.025 & 3.36 \\
Last token   & 0.100 & 0.485 & 3.42 \\
Mean pool    & \textbf{0.105} & \textbf{0.485} & \textbf{3.34} \\
Bidirectional CLS & 0.100 & 0.485 & 3.37 \\
\bottomrule
\end{tabular}
\end{center}

CLS pooling: dead at chance ($0.5\%$). All three fixes: $20\times$ improvement.
Mean pooling is the best strategy: best text quality \emph{and} best contrastive retrieval.

\paragraph{Why audio worked but images didn't.}
Audio classification on the Vast run showed $10$--$20\%$ accuracy (vs $2\%$ chance), despite the same CLS token.
The difference: audio used \texttt{is\_causal=False} (bidirectional), so CLS could see all audio patches.
The bug was specific to text encoding for contrastive alignment.

\subsection{Bug 2: RoPE Position Shift (Secondary)}

The \texttt{encode()} method prepends CLS and MOD tokens before passing through the backbone, shifting all text tokens to RoPE positions $+2$.
Confirmed experimentally: $+0.078$ nats degradation at toy scale.
Three RoPE strategies (naive, neutral, 2D) were tested; all showed identical R@1 because the CLS blindness was the dominant failure.

\subsection{V4 Design: Corrected Multimodal ORN}
\label{sec:v4-design}

The corrected design incorporates lessons from all three multimodal iterations:

\begin{center}\small
\begin{tabular}{lll}
\toprule
\textbf{Component} & \textbf{V3 (broken)} & \textbf{V4 (corrected)} \\
\midrule
Text pooling (contrastive) & CLS pos 0, causal & Mean pool over all positions \\
Image pooling              & CLS pos 0, bidir  & Mean pool over all positions \\
Text LM loss               & Causal attention   & Causal attention (unchanged) \\
RoPE for text              & 1D from pos 0      & 1D from pos 0 (no prefix shift) \\
RoPE for images            & 1D (meaningless)    & Identity (no positional encoding) \\
Input normalisation        & Shared RMSNorm      & Shared RMSNorm (from V2) \\
CLS/MOD prefix tokens      & Prepended           & Removed --- mean pool needs no CLS \\
\bottomrule
\end{tabular}
\end{center}

\paragraph{Key change: no prefix tokens.}
Removing CLS and MOD tokens eliminates both bugs simultaneously.
Without prefix tokens, there is no position shift and no blind CLS.
Modality information enters through the input projection (per-modality linear layers), not through special tokens.
Mean pooling over all content positions provides the contrastive embedding.

\subsection{Proposed Experiments: V4 Validation}

Before committing GPU time, all experiments run on MPS at small scale.

\paragraph{Exp V4-1: Mean pool contrastive with real images (MPS, 2h).}
Train a 12M ORN on TinyStories + COCO patches with paired captions.
Mean pooling for contrastive, causal for text LM.
No CLS/MOD prefix tokens.
\begin{itemize}[nosep]
\item \textbf{Target:} R@1 $> 30\%$ on held-out COCO val. Confirms the fix works with real data.
\item \textbf{Kill criterion:} R@1 $< 10\%$ after 5K steps. If real images don't learn, the problem is deeper.
\end{itemize}

\paragraph{Exp V4-2: Temperature sweep (MPS, 1h).}
Sweep $\tau \in \{0.01, 0.03, 0.07, 0.15, 0.30\}$ on fixed architecture.
\begin{itemize}[nosep]
\item $\tau = 0.07$ was used in V3 without justification. CLIP uses $\tau = 0.01$--$0.03$ at this scale.
\item Measure R@1 and contrastive loss convergence at each temperature.
\end{itemize}

\paragraph{Exp V4-3: Joint training balance (MPS, 3h).}
Sweep text fraction $\in \{0.5, 0.7, 0.8, 0.9\}$.
\begin{itemize}[nosep]
\item V3 used $80\%$ text. But with correct pooling, the contrastive loss may need more budget.
\item Measure: text val loss AND R@1 at each mix. Find the Pareto frontier.
\end{itemize}

\paragraph{Exp V4-4: Image RoPE ablation (MPS, 1h).}
Compare identity RoPE for images vs 2D spatial RoPE vs no RoPE.
\begin{itemize}[nosep]
\item The earlier RoPE experiment was contaminated by the CLS bug. Re-run with mean pooling.
\item Now that contrastive actually works, we can measure whether RoPE strategy affects image quality.
\end{itemize}

\paragraph{Exp V4-5: Modality-specific attention pattern (MPS, 1h).}
Compare:
\begin{enumerate}[nosep, label=(\alph*)]
\item Causal for text, bidirectional for images (V3 approach, but with mean pooling)
\item Bidirectional for all contrastive encoding, causal only for LM loss
\item Causal for everything (images treated as autoregressive patches, raster order)
\end{enumerate}

\paragraph{GPU deployment protocol.}
Only after all five experiments pass on MPS:
\begin{enumerate}[nosep]
\item Verify R@1 $> 20\%$ on real images at 12M scale
\item Verify text val is not degraded vs text-only training
\item Verify no gradient explosions or loss plateaus in first 1K steps
\item Write training script with \textbf{per-modality eval from step 1} (the V3 lesson)
\item Launch on Vast with automatic checkpoint + eval every 500 steps
\end{enumerate}


% ======================================================================
\appendix
\section{Connection to Embeddings and Contrastive Training}
\label{sec:contrastive}

\subsection{The Missing Training Signal}

Next-token prediction teaches $\M$ \emph{within-sequence coupling}: which tokens should attend to which within a single text.
It has zero gradient signal for \emph{between-sequence relationships}: whether two texts are about the same topic, have the same structure, or describe the same image from different modalities.

This explains two persistent failures:
\begin{enumerate}[nosep]
    \item \textbf{Fiber disjointness.} $\M$ discovers that image patches and text tokens couple with $>$99.6\% similarity, but the modalities land in different regions of $\R^d$.
    The coupling is shared; the fibers are disjoint.
    $\M$ never received a signal that an image of a dog and the word ``dog'' should be \emph{near each other}.
    \item \textbf{Weak embeddings.} ORN V2 (108M, 4B tokens FineWeb-Edu) achieves 0.31 on STSBenchmark---below GloVe.
    The model can predict the next word but cannot judge whether two sentences mean the same thing.
\end{enumerate}

Contrastive training provides the missing signal.
It tells $\M$: these two sequences should map to nearby representations (positive pair); these two should map far apart (negative pair).
This is \textbf{between-sequence coupling}---the global analogue of $\M$'s local coupling geometry.

\subsection{Contrastive Training as Part of Pretraining}

The standard approach treats contrastive loss as a post-hoc fine-tuning step: train the LM, then fine-tune for embeddings.
We argue it should be part of ORN pretraining:

\begin{itemize}[nosep]
    \item \textbf{Next-token prediction} teaches $\M$ local structure: syntax, coreference, what-follows-what.
    \item \textbf{Contrastive loss} teaches $\M$ global structure: which sequences are similar, which are different, which modalities describe the same content.
    \item Both signals shape $\M$'s eigenspectrum. Without contrastive signal, the top modes encode communicative form (question vs narrative) but the tail modes fail to cluster topic cleanly. With contrastive signal, the tail modes should develop explicit topic geometry.
\end{itemize}

For multimodal training specifically, contrastive loss is what builds the \emph{on-ramps} between modalities.
The coupling geometry ($\M$) is shared, but the fiber space (where each modality's representations live) remains disjoint without explicit pressure to align.
InfoNCE on cross-modal pairs provides this pressure.

\subsection{Spectral Hard Negatives}

$\M$'s eigenspectrum provides a unique mechanism for hard negative mining---a critical component of embedding quality:

\begin{itemize}[nosep]
    \item Two texts close in top modes (same form) but far in tail modes (different topic) are natural hard negatives for topic retrieval.
    \item Two texts close in tail modes (same topic) but far in top modes (different form) are hard negatives for form retrieval.
    \item Cross-modal pairs close in top modes (shared coupling) but far in tail modes (disjoint fibers) are hard negatives for alignment.
\end{itemize}

No external model (GPT-4) is needed to generate hard negatives---$\M$'s spectral decomposition provides them for free.
This is a direct advantage of the shared coupling architecture over standard transformers.

\subsection{Conjecture: Orbit Isometry as Self-Supervised Contrastive}
\label{sec:orbit-contrastive}

\begin{conjecture}[Self-contrastive ORN]
Two sequences whose orbit trajectories through $\M$'s eigenspectrum are isometric---same eigenmode activation pattern at each depth---are structurally equivalent.
The orbit trajectory itself provides a self-supervised contrastive signal, requiring no external labels or paired data.
\end{conjecture}

\paragraph{The mechanism.}
At each layer $\ell$, project the hidden state $h^{(\ell)}$ onto $\M$'s top-$k$ eigenmodes, yielding an eigenmode activation vector $a^{(\ell)} \in \R^k$.
The orbit trajectory of a sequence is the matrix $\mathcal{T}(x) = [a^{(1)}; a^{(2)}; \ldots; a^{(L)}] \in \R^{L \times k}$.
Two sequences $x, y$ are a positive pair if $\|\mathcal{T}(x) - \mathcal{T}(y)\|$ is small; a negative pair if large.

\paragraph{Connection to CORN.}
The semigroup composition property $T(s{+}t) \approx T(s) \circ T(t)$ means that sequences whose orbits compose in the same way are geometrically related in $\M$'s coupling space.
Orbit isometry implies semigroup equivalence: the two sequences traverse the same path through coupling geometry, activating the same eigenmodes at the same depths.

\paragraph{Why this matters.}
If confirmed, the ORN becomes \textbf{self-contrastive}: the same shared $\M$ that provides within-sequence coupling (via next-token prediction) also provides between-sequence similarity (via orbit isometry).
Next-token prediction crystallises $\M$'s local coupling geometry.
Orbit-similarity crystallises $\M$'s global sequence organisation.
Both arise from one operator.
No external supervision, no paired data, no separate contrastive dataset.

\paragraph{Experimental evidence so far.}
Orbit trajectory embeddings ($\mathcal{T}(x) \in \R^{L \times k}$, flattened) achieve 97.7\% kNN accuracy on communicative form classification---the best of any representation tested, beating static top-mode projection (88.3\%) and full embedding (90.6\%).
The top-5 most similar trajectory pairs are all correct form AND topic matches (cosine similarity ${>}0.99$).

However, orbit trajectory scores 0.247 on STSBenchmark (below mean-pool baseline at 0.314).
STS measures semantic content similarity, not structural equivalence.
This clarifies the role: orbit isometry provides \textbf{training pairs} (structurally similar sequences), not a drop-in embedding replacement.
The pairs must be used with contrastive training to teach $\M$ global semantic organisation.

\paragraph{Proposed ablation.}
Four stages, each measured on STS and mode-selective retrieval:
\begin{enumerate}[nosep]
    \item \textbf{Frozen $\M$ + orbit-derived pairs}: mine positive pairs from trajectory similarity, train readout head only.
    Tests whether the orbit signal is useful for semantic similarity.
    \item \textbf{Frozen $\M$ + SimCSE}: same text, different dropout.
    Baseline contrastive without orbit signal.
    \item \textbf{Frozen $\M$ + orbit + SimCSE}: both signals combined.
    \item \textbf{Unfrozen $\M$ (low LR) + combined}: let $\M$ update from between-sequence feedback.
    Risk: feedback loop ($\M$ changes $\to$ trajectories change $\to$ pairs change $\to$ $\M$ changes).
\end{enumerate}

The partitioning isolates each contribution.
For multimodal, the same ablation applies: orbit isometry across modalities (image trajectory $\approx$ text trajectory) would provide cross-modal positive pairs without paired image-text data.

\paragraph{Risks.}
Orbit trajectories may be too uniform (all sequences converge to $\M$'s dominant eigenmode, giving high similarity everywhere) or too noisy (sequence-specific perturbations dominate, masking structural similarity).
The eigenmode projection may not be discriminative enough at each layer to distinguish structurally equivalent sequences from merely similar ones.
The STSBenchmark result (0.247) shows that structural equivalence $\neq$ semantic similarity---the orbit signal alone is insufficient for content-level judgments.

\end{document}
